% Theories as Regularizers
% Amedeo Andriollo and Juan F. Imbet
% Paris Dauphine University - PSL

\documentclass[12pt]{article}

% ── Page layout ──────────────────────────────────────────────
\usepackage[margin=1in]{geometry}
\usepackage{setspace}
\doublespacing

% ── Mathematics ──────────────────────────────────────────────
\usepackage{amsmath,amssymb,amsthm,mathtools}
\usepackage{bm}

% ── Tables and figures ───────────────────────────────────────
\usepackage{booktabs,tabularx,multirow,array}
\usepackage{graphicx}
\usepackage{caption,subcaption}

% ── Bibliography ─────────────────────────────────────────────
\usepackage[authoryear,round]{natbib}
\bibliographystyle{plainnat}

% ── Cross-references and links ───────────────────────────────
\usepackage{hyperref}
\hypersetup{
  colorlinks=true,
  linkcolor=blue!60!black,
  citecolor=blue!60!black,
  urlcolor=blue!60!black,
  pdfauthor={Amedeo Andriollo and Juan F. Imbet},
  pdftitle={Theories as Regularizers}
}
\usepackage{cleveref}

% ── Miscellaneous ────────────────────────────────────────────
\usepackage{enumerate,enumitem}
\usepackage{xcolor}
\usepackage{dsfont}           % indicator function \mathds{1}

% ── Theorem environments ────────────────────────────────────
\theoremstyle{plain}
\newtheorem{theorem}{Theorem}
\newtheorem{proposition}[theorem]{Proposition}
\newtheorem{lemma}[theorem]{Lemma}
\newtheorem{corollary}[theorem]{Corollary}

\theoremstyle{definition}
\newtheorem{definition}[theorem]{Definition}
\newtheorem{assumption}[theorem]{Assumption}
\newtheorem{example}[theorem]{Example}

\theoremstyle{remark}
\newtheorem{remark}[theorem]{Remark}

% ── Custom commands ──────────────────────────────────────────
\newcommand{\E}{\mathbb{E}}
\newcommand{\R}{\mathbb{R}}
\newcommand{\Hcal}{\mathcal{H}}
\newcommand{\Kcal}{\mathcal{K}}
\newcommand{\Fcal}{\mathcal{F}}
\newcommand{\ind}{\mathds{1}}
\DeclareMathOperator*{\argmin}{arg\,min}
\DeclareMathOperator*{\argmax}{arg\,max}

% ══════════════════════════════════════════════════════════════
%  Title
% ══════════════════════════════════════════════════════════════

\title{\textbf{Theories as Regularizers}\thanks{%
We thank [acknowledgments].
Andriollo: Paris Dauphine University -- PSL, \texttt{amedeo.andriollo@dauphine.psl.eu}.
Imbet: Paris Dauphine University -- PSL, \texttt{juan.imbet@dauphine.psl.eu}.}}
\author{Amedeo Andriollo \and Juan F.\ Imbet}
\date{\today}

\begin{document}

\maketitle

% ── Abstract ─────────────────────────────────────────────────
\begin{abstract}
\begin{singlespace}
\noindent
Machine-learning methods for return prediction ignore economic theory. Structural asset-pricing models encode deep knowledge about expected returns but are individually misspecified. We let the data decide which theories help and how much. We embed over fifty model-implied restrictions---from consumption-based, production-based, intermediary, behavioral, and demand-system frameworks---as soft penalties in a kernel ridge regression estimator. Each restriction carries its own penalty weight, selected by cross-validation: theories that hurt prediction are silenced; theories that help are amplified. The result is the first empirical ranking of economic theories by their marginal out-of-sample forecasting value. Theory-informed penalization delivers higher out-of-sample predictive accuracy and Sharpe ratios than both agnostic machine learning and standard kernel methods. Consumption-based and intermediary restrictions contribute the most. Theory helps disproportionately when data are scarce and during crises---precisely when investors need guidance most.

\medskip
\noindent\textbf{JEL Classification:} G12, G17, C14, C52.

\noindent\textbf{Keywords:} Asset pricing, kernel ridge regression, structural restrictions, machine learning, model selection.
\end{singlespace}
\end{abstract}

\thispagestyle{empty}
\newpage
\setcounter{page}{1}

% ══════════════════════════════════════════════════════════════
%  Main body
% ══════════════════════════════════════════════════════════════

% ══════════════════════════════════════════════════════════════
%  Section 1: Introduction
% ══════════════════════════════════════════════════════════════

\section{Introduction}\label{sec:introduction}

Two traditions compete for the task of predicting asset returns. Statistical machine-learning methods---neural networks, random forests, gradient-boosted trees---estimate flexible functions of observable characteristics and deliver strong out-of-sample (OOS) performance, yet they treat economic theory as irrelevant input \citep{gu2020empirical}. Structural asset-pricing models---CAPM, long-run risk, habit formation, intermediary-based pricing---derive sharp restrictions on the pricing kernel and expected returns, yet each model is misspecified in isolation. One tradition has power without guidance; the other has guidance without robustness. Neither alone is satisfactory.

We propose a framework that draws on both. Our estimator, Theory-Informed Kernel Ridge Regression (Theory-KRR), embeds fifty-six model-implied structural restrictions as soft penalties in the KRR objective. Each restriction $j$ enters with its own penalty multiplier~$\mu_j$, and the entire vector $\bm{\mu} = (\mu_1, \dots, \mu_J)$ is selected by cross-validation. The mechanism is automatic: restrictions from misspecified theories receive~$\mu_j \to 0$---the data shut them down---while restrictions from empirically useful theories receive~$\mu_j > 0$, steering the estimator toward economically sensible regions of the reproducing kernel Hilbert space~$\Hcal$. No single structural model needs to be correct. The data adjudicate.

The restrictions we impose are not generic shape constraints such as monotonicity or convexity \citep{bryzgalova2023asset}. They are genuine model-implied restrictions: Euler-equation moment conditions, factor-loading constraints, demand-system equilibrium conditions, and stochastic discount factor bounds, each derived from a specific economic theory. By cataloging fifty-six such restrictions across consumption-based, production-based, intermediary, behavioral, and demand-system frameworks, we construct a comprehensive menu of theoretical guidance. Cross-validation then selects which theories---and which specific restrictions within each theory---contribute to OOS prediction.

A methodological parallel clarifies the logic. The LLM-Lasso of \citet{bybee2025llmlasso} replaces the agnostic $\ell_1$ penalty in Lasso with heterogeneous, text-informed penalties: large-language-model predictions set a prior over which covariates matter, and cross-validation adjusts a single scalar that scales the text-derived penalty vector. Our approach applies the same principle in a different domain. Where LLM-Lasso uses text priors to inform penalization in \emph{parameter space}, Theory-KRR uses economic-theory priors to inform penalization in \emph{function space}. The mathematical structure is analogous: both replace a single, uninformed regularization parameter with a vector of informed, heterogeneous penalties, and both let cross-validation arbitrate between the prior and the data. Theory-KRR is, in this sense, the function-space generalization of LLM-Lasso.

What does this framework deliver? Theory-KRR outperforms standard KRR and several benchmark ML methods in OOS return prediction, measured by $R^2_{\text{OOS}}$ and portfolio Sharpe ratios. The improvement is not uniform: theory-informed penalization helps most in small samples and during financial crises, precisely when purely statistical methods overfit and economic structure is most valuable. Beyond predictive gains, the cross-validated multipliers~$\bm{\mu}$ produce an empirical ranking of economic theories by their marginal OOS forecasting value. Consumption-based restrictions \citep{campbellcochrane1999,bansalyaron2004} and intermediary-based restrictions \citep{hekrishnamurthy2013,adrianetulamuir2014} receive the largest multipliers, indicating that these theories provide the most useful guidance for return prediction. Production-based and behavioral restrictions contribute modestly. Several models receive~$\mu_j$ indistinguishable from zero, consistent with those theories being empirically irrelevant for forecasting in our sample.

This paper makes three contributions. First, we develop a unified framework that embeds fifty-six structural restrictions from diverse asset-pricing models as soft penalties in a single KRR estimator. The framework nests standard KRR (all~$\mu_j = 0$) and hard-constrained estimation (any~$\mu_j \to \infty$) as special cases, and it accommodates restrictions of heterogeneous functional form---moment conditions, inequality constraints, and functional equations---within a single objective. Second, we establish a formal connection between theory-informed function-space penalization and the LLM-Lasso approach to parameter-space penalization. This connection unifies two strands of the literature---structural econometrics and text-based ML---under a common principle of informed regularization. Third, we produce the first empirical ranking of economic theories by their marginal contribution to OOS forecasting, providing a new lens through which to evaluate the empirical content of competing asset-pricing models.

Our work speaks to several strands of the literature. The ML asset-pricing literature, initiated by \citet{gu2020empirical}, demonstrates that flexible nonparametric methods substantially outperform linear models in predicting the cross-section of returns. Subsequent work explores neural-network architectures \citep{chen2024deep}, autoencoder factor models with no-arbitrage constraints \citep{Gu2021autoencoder}, conditional factor models \citep{kelly2019characteristics}, and Gaussian process regression applied to equity returns \citep{Filipovic2022gaussian}. \citet{Kelly2024virtue} provide a theoretical foundation for this success, showing that complex, overparameterized models are not penalized by benign overfitting and systematically outperform simpler ones. These methods are powerful but agnostic about economic structure. Our contribution is to show that injecting economic theory into the regularization of a kernel method yields both predictive gains and economic interpretability.

A growing body of work recognizes that theory and statistics are complements, not substitutes, for return prediction. The paper closest to ours in spirit is \citet{Chen2026teaching}, who show that pre-training neural networks on synthetic data generated by structural models---and then fine-tuning on empirical data---substantially improves option-pricing accuracy, even when the structural model is misspecified. Our approach shares the same fundamental insight---misspecified theory still helps---but differs in three important respects. First, we use soft penalization in the objective function rather than pre-training: structural content enters directly as penalty terms with cross-validated weights, not as an initialization from synthetic data. Second, we target the cross-section of equity returns rather than option pricing. Third, the convexity of our KRR objective guarantees a global optimum and yields tractable asymptotics, whereas neural-network pre-training is inherently non-convex. Most importantly, our penalty weights provide a continuous measure of each theory's marginal contribution to prediction, whereas transfer learning offers a binary use-or-not decision for each structural model.

A second closely related paper is \citet{Bryzgalova2023bayesian}, who develop a Bayesian framework that averages over 2.25~quadrillion linear asset-pricing models formed by subsets of 51~observable factors. Their identification-strength-based priors impose heterogeneous shrinkage---factors that are weakly identified receive stronger shrinkage---and the resulting Bayesian model-averaging stochastic discount factor dominates individual factor models out of sample. Like us, they produce an empirical ranking of economic models. The key differences are the objects being ranked and the estimation framework. Bryzgalova, Huang, and Julliard rank observable factor portfolios within a linear SDF; we rank structural restrictions from named theoretical models within a nonparametric RKHS. Their priors are calibrated from factor identification strength; our penalty weights are selected by temporal cross-validation, tying the ranking directly to out-of-sample predictive performance.

More broadly, the question of which factors or models survive in a high-dimensional setting has been explored from several angles. \citet{Kozak2020shrinking} propose economically motivated shrinkage of SDF coefficients proportional to the eigenvalues of characteristic-based principal components, motivated by no-arbitrage bounds on Sharpe ratios. \citet{Feng2020taming} develop a double-selection LASSO procedure to test whether newly proposed factors have marginal explanatory power beyond existing ones. \citet{Harvey2016cross} document severe data snooping in factor discovery and argue for higher significance thresholds. \citet{Barillas2018comparing} and \citet{Kan2013pricing} provide formal Bayesian and frequentist frameworks for pairwise model comparison. Our paper complements this literature by moving from factor-level selection in linear models to theory-level selection in a nonparametric estimator, where the ``theories'' are entire families of structural restrictions rather than individual tradable factors.

The econometrics literature on constrained nonparametric estimation provides our mathematical foundation. \citet{babii2024structural} develop the theory of penalized kernel regression with functional constraints, establishing convergence rates and conditions under which constraints improve estimation. \citet{gabaix2024asset} show how demand-system equilibrium conditions can be imposed on asset-pricing estimators, and \citet{bryzgalova2023asset} impose shape restrictions (monotonicity, convexity) on pricing-kernel estimators. In the penalized GMM literature, \citet{Caner2009lasso} and \citet{Caner2014adaptive} develop LASSO-type and elastic-net GMM estimators with oracle properties for moment selection. \citet{Freyberger2020dissecting} use adaptive group LASSO to select characteristics nonparametrically. Our work extends this line by moving from reduced-form shape restrictions to model-implied structural restrictions, and by allowing the data to select among competing restrictions via cross-validation rather than imposing all restrictions with equal weight. The representer theorem, generalized to our multi-penalty setting following \citet{Scholkopf2002learning} and the foundational results of \citet{Kimeldorf1971correspondence} and \citet{Wahba1990spline}, ensures that the solution lies in the span of kernel evaluations even with the added structural penalties.

The paper proceeds as follows. Section~\ref{sec:framework} presents the mathematical framework, defining the Theory-KRR objective and establishing its properties. Section~\ref{sec:restrictions} catalogs the fifty-six structural restrictions we draw from the asset-pricing literature, organized by model class. Section~\ref{sec:estimation} addresses practical estimation, including cross-validation of the multiplier vector~$\bm{\mu}$, computational strategies, and implementation details. Section~\ref{sec:empirical} describes the data and baseline models. Section~\ref{sec:results} presents the main empirical results: OOS predictive performance, the theory ranking, and robustness exercises. Section~\ref{sec:conclusion} concludes.

%!TEX root = ../main.tex
%----------------------------------------------------------------------
%  Section: The Theory-Informed KRR Estimator
%  Paper:   Theory-Informed Kernel Ridge Regression for Asset Pricing
%  Authors: Amedeo Andriollo and Juan F. Imbet
%----------------------------------------------------------------------

\section{The Theory-Informed Kernel Ridge Regression Estimator}
\label{sec:framework}

This section develops the econometric framework.
Section~\ref{sub:setup} introduces notation and the prediction environment.
Section~\ref{sub:estimator} defines the theory-informed KRR objective.
Sections~\ref{sub:penalties}--\ref{sub:finite_dim} formalize the structural penalty types
and establish the representer theorem and finite-dimensional reduction that
make the estimator computationally tractable.
Section~\ref{sub:nesting} shows that the framework nests several influential
estimators as special cases.
Sections~\ref{sub:cv}--\ref{sub:inference} address implementation, misspecification
robustness, and inference.

%======================================================================
\subsection{Setup and Notation}
\label{sub:setup}
%======================================================================

We observe a panel of $N$ assets over $T$ periods.
Let $R_{i,t+1}$ denote the excess return of asset $i$ realized at date $t+1$,
and let $\mathbf{X}_{i,t} \in \mathbb{R}^p$ collect firm-level characteristics
and macroeconomic state variables observable at date $t$.
The total sample size is $n = NT$.

\begin{assumption}[Data-generating process]
\label{ass:dgp}
There exists a measurable function $f^{*}: \mathbb{R}^{p} \to \mathbb{R}$
such that
\begin{equation}
\label{eq:dgp}
R_{i,t+1} = f^{*}(\mathbf{X}_{i,t}) + \varepsilon_{i,t+1},
\qquad
\mathbb{E}[\varepsilon_{i,t+1} \mid \mathbf{X}_{i,t}] = 0,
\end{equation}
where $f^{*}(\mathbf{X}_{i,t}) = \mathbb{E}[R_{i,t+1} \mid \mathbf{X}_{i,t}]$
is the conditional expected excess return.
\end{assumption}

We approximate $f^{*}$ within a reproducing kernel Hilbert space (RKHS).

\begin{definition}[RKHS]
\label{def:rkhs}
Let $k: \mathbb{R}^p \times \mathbb{R}^p \to \mathbb{R}$ be a symmetric,
positive-definite kernel.
The RKHS $\mathcal{H}$ associated with $k$ is the completion of
$\operatorname{span}\{k(\mathbf{x}, \cdot) : \mathbf{x} \in \mathbb{R}^p\}$
under the inner product satisfying the reproducing property
$f(\mathbf{x}) = \langle f, k(\mathbf{x}, \cdot) \rangle_{\mathcal{H}}$
for all $f \in \mathcal{H}$.
The RKHS norm $\|f\|_{\mathcal{H}}^2 = \langle f, f \rangle_{\mathcal{H}}$
measures the complexity of $f$.
\end{definition}

We stack observations using the index $s = (i,t)$, so that
$\mathbf{X}_s \equiv \mathbf{X}_{i,t}$ and $R_s \equiv R_{i,t+1}$ for
$s = 1, \ldots, n$.
The $n \times n$ kernel (Gram) matrix is
$\mathbf{K}$, with $[\mathbf{K}]_{s,s'} = k(\mathbf{X}_s, \mathbf{X}_{s'})$.
The return vector is $\mathbf{R} = (R_1, \ldots, R_n)^\top \in \mathbb{R}^n$.

%======================================================================
\subsection{The Estimator}
\label{sub:estimator}
%======================================================================

Our estimator augments the standard KRR objective with a collection of
structural penalty terms, each encoding restrictions implied by a specific
economic model.

\begin{definition}[Theory-Informed KRR]
\label{def:tikrr}
Let $\{C_j\}_{j=1}^{J}$ be a collection of structural penalty functionals
$C_j: \mathcal{H} \to \mathbb{R}_+$,
where each $C_j(f) \geq 0$ with equality when $f$ satisfies the
restrictions of economic model~$j$.
Let $\lambda_{\mathrm{stat}} > 0$ and $\mu_j \geq 0$, $j = 1, \ldots, J$.
The \emph{theory-informed KRR} estimator is
\begin{equation}
\label{eq:tikrr}
\hat{f} \;=\; \arg\min_{f \in \mathcal{H}} \;
\underbrace{\frac{1}{n}\sum_{s=1}^{n}
  \bigl(R_s - f(\mathbf{X}_s)\bigr)^2}_{\text{prediction loss}}
\;+\;
\underbrace{\lambda_{\mathrm{stat}}\,\|f\|_{\mathcal{H}}^2}_{\text{statistical regularization}}
\;+\;
\underbrace{\sum_{j=1}^{J} \mu_j\, C_j(f)}_{\text{structural penalties}}.
\end{equation}
\end{definition}

The first two terms constitute standard kernel ridge regression:
the squared-error prediction loss trades off against the RKHS-norm penalty,
which controls the smoothness of $f$.
The third term is new.
Each $C_j$ penalizes deviations from a specific economic model's
predictions, and the associated multiplier $\mu_j$ governs how strongly
that model's restrictions influence the estimate.
When $\mu_j = 0$, model $j$ is ignored; as $\mu_j \to \infty$,
the restrictions of model $j$ are imposed exactly.

\begin{remark}[Interpretation of multipliers]
\label{rem:multipliers}
The vector $(\lambda_{\mathrm{stat}}, \mu_1, \ldots, \mu_J)$ balances
agnostic statistical regularization against theory-informed regularization.
A large $\lambda_{\mathrm{stat}}$ relative to $\mu_j$ yields an estimate
close to standard KRR; large $\mu_j$ relative to $\lambda_{\mathrm{stat}}$
forces the estimator to respect model $j$'s restrictions even if doing
so worsens in-sample fit.
Cross-validation selects the balance that maximizes out-of-sample
predictive accuracy (Section~\ref{sub:cv}).
\end{remark}

%======================================================================
\subsection{Structural Penalty Types}
\label{sub:penalties}
%======================================================================

We formalize four classes of structural penalties.
Each encodes restrictions from a distinct family of asset-pricing models.
A key requirement---verified explicitly below---is that every $C_j(f)$
depends on $f$ only through its evaluations at the observed data points
$\{f(\mathbf{X}_s)\}_{s=1}^n$.

%----------------------------------------------------------------------
\subsubsection{Type~A: Euler Equation Penalties}
\label{sub:euler}
%----------------------------------------------------------------------

Let model $j$ specify a stochastic discount factor (SDF)
$M_{j,t+1} = M_j(\theta_j, \mathbf{X}_t, \mathbf{X}_{t+1})$
parameterized by $\theta_j$, with a preliminary estimator
$\hat{\theta}_j$ obtained by GMM or maximum likelihood.
The Euler equation requires
$\mathbb{E}_t[M_{j,t+1}\, R_{i,t+1}^{\mathrm{gross}}] = 1$
for every asset $i$ and date $t$,
where $R_{i,t+1}^{\mathrm{gross}} = f(\mathbf{X}_{i,t}) + R_t^f$
is the gross return and $R_t^f$ is the risk-free rate.

\begin{definition}[Euler equation penalty]
\label{def:euler_penalty}
For SDF model $j$ with pre-estimated parameter $\hat{\theta}_j$, define
\begin{equation}
\label{eq:euler_penalty}
C_j^{\mathrm{Euler}}(f) \;=\;
\sum_{t=1}^{T}
\left\|
\frac{1}{N}\sum_{i=1}^{N}
M_j\bigl(\hat{\theta}_j, \mathbf{X}_t, \mathbf{X}_{t+1}\bigr)
\cdot
\bigl(f(\mathbf{X}_{i,t}) + R_t^f\bigr) - 1
\right\|^2.
\end{equation}
\end{definition}

\noindent
The penalty is zero precisely when the cross-sectional Euler equation
holds at every date $t$ given $\hat{\theta}_j$.
Since $M_j(\hat{\theta}_j, \mathbf{X}_t, \mathbf{X}_{t+1})$ and $R_t^f$ are data,
the penalty depends on $f$ only through $\{f(\mathbf{X}_{i,t})\}_{i,t}$.

\begin{remark}[Coverage]
Type~A penalties accommodate consumption-based models (CCAPM, long-run risk,
rare disasters, habit formation), as well as production-based SDFs derived
from firms' intertemporal optimality conditions.
Each distinct SDF specification generates a separate penalty $C_j^{\mathrm{Euler}}$.
\end{remark}

%----------------------------------------------------------------------
\subsubsection{Type~B: Production and Investment FOC Penalties}
\label{sub:production}
%----------------------------------------------------------------------

Production-based asset pricing relates expected returns to observable
investment and profitability variables through firms' first-order conditions.
Let model $j$ imply a mapping $g_j: \mathbb{R}^{d_j} \to \mathbb{R}$
from investment characteristics $\mathbf{Z}_{i,t}^{(j)}$
(e.g., investment-to-capital ratio $I_{i,t}/K_{i,t}$, return on equity
$\mathrm{ROE}_{i,t}$, asset growth) to expected returns.

\begin{definition}[Production FOC penalty]
\label{def:prod_penalty}
For production model $j$ with predicted returns $g_j(\mathbf{Z}_{i,t}^{(j)})$, define
\begin{equation}
\label{eq:prod_penalty}
C_j^{\mathrm{prod}}(f) \;=\;
\frac{1}{n}\sum_{s=1}^{n}
\bigl(
f(\mathbf{X}_s) - g_j\bigl(\mathbf{Z}_s^{(j)}\bigr)
\bigr)^2.
\end{equation}
\end{definition}

\noindent
The penalty is a mean-squared deviation between the nonparametric estimate
and the production model's implied expected return.
Since $g_j(\mathbf{Z}_s^{(j)})$ is computed from data, the penalty depends on $f$
only through $\{f(\mathbf{X}_s)\}_{s=1}^n$.

\begin{remark}[Nesting the $q$-theory framework]
When $g_j$ encodes the first-order condition from the investment CAPM
or the $q$-factor model,
Type~B penalties embed the production approach
to asset pricing directly into the nonparametric objective.
\end{remark}

%----------------------------------------------------------------------
\subsubsection{Type~C: Factor Structure Penalties}
\label{sub:factor}
%----------------------------------------------------------------------

Factor models impose that expected returns obey a linear pricing
relation with respect to exposure to priced factors.
Let model $j$ specify a factor $F_{j,t}$ (e.g., an intermediary capital
ratio factor, market excess return, or a latent factor).
Let $\beta_{i,j}$ denote asset $i$'s loading on $F_{j,t}$,
pre-estimated by time-series regression, and let $\alpha_{i,j}$
be the corresponding intercept.

\begin{definition}[Factor structure penalty]
\label{def:factor_penalty}
For factor model $j$ with estimated loadings
$\hat{\beta}_{i,j}$, $\hat{\alpha}_{i,j}$ ($i=1,\ldots,N$)
and factor realizations $F_{j,t}$ ($t=1,\ldots,T$), define
\begin{equation}
\label{eq:factor_penalty}
C_j^{\mathrm{factor}}(f) \;=\;
\frac{1}{n}\sum_{i=1}^{N}\sum_{t=1}^{T}
\bigl(
f(\mathbf{X}_{i,t}) - \hat{\beta}_{i,j}\, F_{j,t} - \hat{\alpha}_{i,j}
\bigr)^2.
\end{equation}
\end{definition}

\noindent
When asset pricing restricts $\hat{\alpha}_{i,j} = 0$ for all $i$
(exact factor pricing), the penalty reduces to
$\frac{1}{n}\sum_{i,t}(f(\mathbf{X}_{i,t}) - \hat{\beta}_{i,j}\, F_{j,t})^2$,
penalizing deviations from the factor model's predictions of expected returns.
Again, the penalty depends on $f$ only through $\{f(\mathbf{X}_{i,t})\}_{i,t}$.

\begin{remark}[Intermediary asset pricing]
When $F_{j,t}$ is the intermediary capital ratio or the intermediary leverage factor,
Type~C penalties embed intermediary-based
models into the nonparametric estimator.
This is particularly valuable when the intermediary model captures
time-variation in risk premia that firm-level characteristics alone cannot.
\end{remark}

%----------------------------------------------------------------------
\subsubsection{Type~D: Demand-System Equilibrium Penalties}
\label{sub:demand}
%----------------------------------------------------------------------

Demand-system approaches to asset pricing derive equilibrium prices
from investors' portfolio demand functions.
Let $D_i: \mathbb{R}^N \to \mathbb{R}^N$ map the vector of expected returns to investor~$i$'s
optimal portfolio weights, and let $w_i > 0$ denote
investor~$i$'s wealth share, with the index now running over investors
rather than assets.
Let $\mathbf{S}_t \in \mathbb{R}^N$ denote the vector of asset market-capitalization
shares (the supply side) at date $t$.

\begin{definition}[Demand-system equilibrium penalty]
\label{def:demand_penalty}
For demand-system model $j$ with demand functions
$\{D_i^{(j)}\}$ and observed supplies $\{\mathbf{S}_t\}$, let
$\hat{\mathbf{f}}_t = (f(\mathbf{X}_{1,t}), \ldots, f(\mathbf{X}_{N,t}))^\top$
collect expected-return evaluations.  Define
\begin{equation}
\label{eq:demand_penalty}
C_j^{\mathrm{demand}}(f) \;=\;
\sum_{t=1}^{T}
\left\|
\sum_{i} w_i^{(j)}\, D_i^{(j)}\bigl(\hat{\mathbf{f}}_t\bigr) - \mathbf{S}_t
\right\|^2.
\end{equation}
\end{definition}

\noindent
The penalty measures the squared market-clearing failure:
it is zero when the demand system clears at the predicted expected returns.
It depends on $f$ only through $\{f(\mathbf{X}_{i,t})\}_{i,t}$.

\begin{remark}[Connections to the demand-system literature]
The demand-system approach models institutional
investors' demand as a function of characteristics and expected returns.
Type~D penalties encode the equilibrium restriction that aggregate demand
equals supply, translating the demand-system estimator's logic into a
penalty within our nonparametric framework.
When $D_i^{(j)}$ is linear in expected returns, $C_j^{\mathrm{demand}}$
is quadratic in $\hat{\mathbf{f}}_t$ and the overall problem remains convex.
\end{remark}

%======================================================================
\subsection{The Representer Theorem}
\label{sub:representer}
%======================================================================

The following result establishes that the infinite-dimensional
optimization problem~\eqref{eq:tikrr} admits a finite-dimensional
solution, despite the presence of the structural penalties.

\begin{proposition}[Representer theorem for theory-informed KRR]
\label{prop:representer}
Let $\mathcal{H}$ be an RKHS with kernel $k$.
Suppose each structural penalty $C_j: \mathcal{H} \to \mathbb{R}_+$
depends on $f$ only through the vector of evaluations
$(f(\mathbf{X}_1), \ldots, f(\mathbf{X}_n))^\top \in \mathbb{R}^n$.
That is, there exist functions $\tilde{C}_j: \mathbb{R}^n \to \mathbb{R}_+$
such that
\begin{equation}
\label{eq:penalty_eval}
C_j(f) = \tilde{C}_j\bigl(f(\mathbf{X}_1), \ldots, f(\mathbf{X}_n)\bigr)
\qquad \text{for all } f \in \mathcal{H}, \; j = 1, \ldots, J.
\end{equation}
Then any minimizer of~\eqref{eq:tikrr} admits the representation
\begin{equation}
\label{eq:representer}
\hat{f}(\cdot) = \sum_{s=1}^{n} \hat{\alpha}_s\, k(\mathbf{X}_s, \cdot),
\end{equation}
for some $\hat{\boldsymbol{\alpha}} = (\hat{\alpha}_1, \ldots, \hat{\alpha}_n)^\top
\in \mathbb{R}^n$.
\end{proposition}

\begin{proof}
\emph{The proof is given in Appendix~\ref{app:proofs}.}
\end{proof}

\begin{remark}[Verification for Types~A--D]
\label{rem:verify_representer}
The condition~\eqref{eq:penalty_eval} holds for all four penalty types:
\begin{enumerate}[label=(\alph*)]
\item \emph{Type~A (Euler):}
  $C_j^{\mathrm{Euler}}(f)$ depends on $f$ through
  $\{f(\mathbf{X}_{i,t})\}_{i,t}$, since
  $M_j(\hat{\theta}_j, \mathbf{X}_t, \mathbf{X}_{t+1})$ and $R_t^f$
  are data.
\item \emph{Type~B (Production):}
  $C_j^{\mathrm{prod}}(f)$ is explicitly a function of
  $\{f(\mathbf{X}_s)\}_{s}$ and data $\{g_j(\mathbf{Z}_s^{(j)})\}_s$.
\item \emph{Type~C (Factor):}
  $C_j^{\mathrm{factor}}(f)$ involves $\{f(\mathbf{X}_{i,t})\}_{i,t}$
  and pre-estimated $\hat{\beta}_{i,j}$, $\hat{\alpha}_{i,j}$, $F_{j,t}$.
\item \emph{Type~D (Demand):}
  $C_j^{\mathrm{demand}}(f)$ depends on $f$ through
  $\hat{\mathbf{f}}_t = (f(\mathbf{X}_{1,t}), \ldots, f(\mathbf{X}_{N,t}))^\top$.
\end{enumerate}
In each case, the penalty is a function of the $n$-dimensional evaluation
vector, so Proposition~\ref{prop:representer} applies.
This extends classical representer theorems to accommodate theory-informed
structural penalties.
\end{remark}

%======================================================================
\subsection{Finite-Dimensional Reduction}
\label{sub:finite_dim}
%======================================================================

By Proposition~\ref{prop:representer}, it suffices to optimize
over $\boldsymbol{\alpha} \in \mathbb{R}^n$.
Substituting $f(\cdot) = \sum_s \alpha_s\, k(\mathbf{X}_s, \cdot)$ into
the objective yields the following.

\begin{proposition}[Finite-dimensional problem]
\label{prop:finite_dim}
Under the representer theorem (Proposition~\ref{prop:representer}),
the theory-informed KRR estimator~\eqref{eq:tikrr} reduces to
\begin{equation}
\label{eq:finite_dim}
\hat{\boldsymbol{\alpha}} \;=\; \arg\min_{\boldsymbol{\alpha} \in \mathbb{R}^n} \;
\frac{1}{n}\|\mathbf{R} - \mathbf{K}\boldsymbol{\alpha}\|^2
\;+\;
\lambda_{\mathrm{stat}}\,\boldsymbol{\alpha}^\top \mathbf{K}\, \boldsymbol{\alpha}
\;+\;
\sum_{j=1}^{J} \mu_j\, \tilde{C}_j(\mathbf{K}\boldsymbol{\alpha}),
\end{equation}
where $\mathbf{K}$ is the $n \times n$ kernel matrix,
$\hat{\mathbf{f}} = \mathbf{K}\boldsymbol{\alpha}$ is the $n \times 1$
vector of fitted values, and
$\tilde{C}_j(\hat{\mathbf{f}}) = C_j(f)$ expresses each penalty
in terms of fitted values.
\end{proposition}

\begin{proof}
\emph{The proof is given in Appendix~\ref{app:proofs}.}
\end{proof}

%----------------------------------------------------------------------
\subsubsection{Convexity and First-Order Conditions for Quadratic Penalties}
\label{sub:foc}
%----------------------------------------------------------------------

When all active structural penalties are quadratic in the fitted values
$\hat{\mathbf{f}} = \mathbf{K}\boldsymbol{\alpha}$, the finite-dimensional
problem~\eqref{eq:finite_dim} becomes a convex quadratic program with
a closed-form solution.

For Type~A penalties, define the $T \times n$ matrix $\mathbf{A}_j$
with rows indexed by $t$:
\begin{equation}
\label{eq:A_matrix}
[\mathbf{A}_j]_{t,s} =
\begin{cases}
\dfrac{1}{N}\, M_j\bigl(\hat{\theta}_j, \mathbf{X}_t, \mathbf{X}_{t+1}\bigr)
& \text{if } s = (i,t) \text{ for some } i, \\[4pt]
0 & \text{otherwise},
\end{cases}
\end{equation}
and the $T \times 1$ vector $\mathbf{b}_j$ with
$[\mathbf{b}_j]_t = 1 - \frac{1}{N}\sum_i M_j(\hat{\theta}_j,
\mathbf{X}_t, \mathbf{X}_{t+1})\, R_t^f$.
Then
\begin{equation}
\label{eq:euler_quad}
C_j^{\mathrm{Euler}}(f) =
\|\mathbf{A}_j\, \hat{\mathbf{f}} - \mathbf{b}_j\|^2
= (\mathbf{A}_j\, \mathbf{K}\boldsymbol{\alpha} - \mathbf{b}_j)^\top
  (\mathbf{A}_j\, \mathbf{K}\boldsymbol{\alpha} - \mathbf{b}_j).
\end{equation}

For Type~B penalties, let $\mathbf{g}_j \in \mathbb{R}^n$ with
$[\mathbf{g}_j]_s = g_j(\mathbf{Z}_s^{(j)})$.  Then
\begin{equation}
\label{eq:prod_quad}
C_j^{\mathrm{prod}}(f) =
\frac{1}{n}\|\mathbf{K}\boldsymbol{\alpha} - \mathbf{g}_j\|^2.
\end{equation}

For Type~C penalties, let $\mathbf{h}_j \in \mathbb{R}^n$ with
$[\mathbf{h}_j]_s = \hat{\beta}_{i,j}\, F_{j,t} + \hat{\alpha}_{i,j}$
where $s = (i,t)$.  Then
\begin{equation}
\label{eq:factor_quad}
C_j^{\mathrm{factor}}(f) =
\frac{1}{n}\|\mathbf{K}\boldsymbol{\alpha} - \mathbf{h}_j\|^2.
\end{equation}

\begin{proposition}[First-order conditions]
\label{prop:foc}
Suppose every active penalty ($\mu_j > 0$) is of Type~A, B, or C
(or more generally, quadratic in $\hat{\mathbf{f}}$).
Define the aggregated penalty matrix
\begin{equation}
\label{eq:Omega}
\boldsymbol{\Omega} \;=\;
\sum_{j \in \mathcal{J}^A} \mu_j\, \mathbf{A}_j^\top \mathbf{A}_j
\;+\;
\sum_{j \in \mathcal{J}^B} \frac{\mu_j}{n}\, \mathbf{I}_n
\;+\;
\sum_{j \in \mathcal{J}^C} \frac{\mu_j}{n}\, \mathbf{I}_n,
\end{equation}
where $\mathcal{J}^A$, $\mathcal{J}^B$, $\mathcal{J}^C$ index
Type~A, B, and C penalties respectively, and the aggregated target vector
\begin{equation}
\label{eq:omega_vec}
\boldsymbol{\omega} \;=\;
\sum_{j \in \mathcal{J}^A} \mu_j\, \mathbf{A}_j^\top \mathbf{b}_j
\;+\;
\sum_{j \in \mathcal{J}^B} \frac{\mu_j}{n}\, \mathbf{g}_j
\;+\;
\sum_{j \in \mathcal{J}^C} \frac{\mu_j}{n}\, \mathbf{h}_j.
\end{equation}
Then the first-order conditions for~\eqref{eq:finite_dim} are
\begin{equation}
\label{eq:foc}
\left[
\frac{1}{n}\,\mathbf{K}
+ \lambda_{\mathrm{stat}}\,\mathbf{I}_n
+ \boldsymbol{\Omega}\,\mathbf{K}
\right] \boldsymbol{\alpha}
\;=\;
\frac{1}{n}\,\mathbf{R}
+ \boldsymbol{\omega},
\end{equation}
and when $\mathbf{K}$ is strictly positive definite, the unique solution is
\begin{equation}
\label{eq:alpha_star}
\hat{\boldsymbol{\alpha}} \;=\;
\left[
\frac{1}{n}\,\mathbf{K}
+ \lambda_{\mathrm{stat}}\,\mathbf{I}_n
+ \boldsymbol{\Omega}\,\mathbf{K}
\right]^{-1}
\left(
\frac{1}{n}\,\mathbf{R} + \boldsymbol{\omega}
\right).
\end{equation}
\end{proposition}

\begin{proof}
\emph{The proof is given in Appendix~\ref{app:proofs}.}
\end{proof}

\begin{remark}[Computational cost]
\label{rem:computation}
The dominant cost is solving the $n \times n$ linear system~\eqref{eq:alpha_star},
which requires $O(n^3)$ operations.
When $n$ is large, one can exploit the eigendecomposition
$\mathbf{K} = \mathbf{U}\boldsymbol{\Lambda}\mathbf{U}^\top$
(computed once at cost $O(n^3)$) and solve in the rotated basis,
where the matrix becomes diagonal plus a low-rank perturbation from
the Type~A penalties.
Alternatively, random Fourier features or Nystr\"{o}m approximations
reduce the effective dimension.
\end{remark}

\begin{remark}[Type~D penalties and convexity]
\label{rem:type_d_convexity}
When demand functions $D_i^{(j)}$ are linear in expected returns,
$C_j^{\mathrm{demand}}$ is quadratic in $\hat{\mathbf{f}}$ and the
problem remains a convex QP with a structure analogous to~\eqref{eq:foc}.
For nonlinear demand systems, the problem is generally nonconvex and
requires iterative methods (e.g., projected gradient descent or
alternating minimization), though convexity of the remaining terms
provides a useful warm start.
\end{remark}

%======================================================================
\subsection{Nesting Structure}
\label{sub:nesting}
%======================================================================

A key advantage of the framework is that it nests several influential
estimators as special or limiting cases.
This both clarifies the relationship between existing methods and
suggests that the theory-informed KRR estimator can weakly dominate
each special case.

\begin{proposition}[Nesting]
\label{prop:nesting}
The theory-informed KRR estimator~\eqref{eq:tikrr} nests the following:
\begin{enumerate}[label=(\roman*)]
\item \textbf{Standard KRR.}
  Setting $\mu_j = 0$ for all $j = 1, \ldots, J$ yields
  \[
  \hat{f}^{\mathrm{KRR}} = \arg\min_{f \in \mathcal{H}}
  \frac{1}{n}\sum_s (R_s - f(\mathbf{X}_s))^2
  + \lambda_{\mathrm{stat}}\|f\|_{\mathcal{H}}^2,
  \]
  the standard kernel ridge regression estimator.

\item \textbf{Constrained minimum-distance / GMM.}
  Setting $\lambda_{\mathrm{stat}} = 0$ and activating a single
  Type~A penalty with $\mu_j \to \infty$ yields
  \[
  \hat{f}^{\mathrm{GMM}} = \arg\min_{f \in \mathcal{H}}
  \frac{1}{n}\sum_s (R_s - f(\mathbf{X}_s))^2
  \quad \text{s.t.} \quad
  C_j^{\mathrm{Euler}}(f) = 0,
  \]
  which is a nonparametric generalization of GMM-based SDF estimation.
  When $\mathcal{H}$ is restricted to parametric functions,
  this recovers classical GMM.

\item \textbf{Approximate instrumented PCA.}
  Consider a linear kernel $k(\mathbf{x}, \mathbf{x}') = \mathbf{x}^\top \mathbf{x}'$
  so that $f(\mathbf{X}_{i,t}) = \boldsymbol{\beta}^\top \mathbf{X}_{i,t}$
  for some $\boldsymbol{\beta} \in \mathbb{R}^p$, and a Type~C penalty
  where the factor $F_{j,t}$ is a target-PCA factor.
  Then~\eqref{eq:tikrr} becomes
  \begin{equation}
  \label{eq:ipca_nest}
  \hat{\boldsymbol{\beta}} = \arg\min_{\boldsymbol{\beta}}
  \frac{1}{n}\sum_{i,t}(R_{i,t+1} - \boldsymbol{\beta}^\top \mathbf{X}_{i,t})^2
  + \lambda_{\mathrm{stat}}\|\boldsymbol{\beta}\|^2
  + \mu_j \frac{1}{n}\sum_{i,t}
    (\boldsymbol{\beta}^\top \mathbf{X}_{i,t}
     - \hat{\beta}_{i,j} F_{j,t})^2,
  \end{equation}
  a ridge regression with an additional penalty encouraging alignment
  with the target-PCA factor structure.
  This approximates the cross-sectional
  and time-series restrictions in instrumented PCA methods.

\item \textbf{Text-informed penalized regression.}
  Replacing $\mathcal{H}$ with a finite-dimensional linear space
  $\mathbb{R}^p$, substituting the RKHS-norm penalty with an
  $\ell_1$ norm $\|\boldsymbol{\beta}\|_1$, and setting
  $C_j(\boldsymbol{\beta}) = \sum_{k=1}^p \gamma_{j,k}\, |\beta_k|$
  where $\gamma_{j,k}$ are text-derived importance weights, yields
  \begin{equation}
  \label{eq:llm_lasso_nest}
  \hat{\boldsymbol{\beta}} = \arg\min_{\boldsymbol{\beta}}
  \frac{1}{n}\sum_s (R_s - \boldsymbol{\beta}^\top \mathbf{X}_s)^2
  + \lambda_{\mathrm{stat}} \|\boldsymbol{\beta}\|_1
  + \mu_j \sum_k \gamma_{j,k}\, |\beta_k|,
  \end{equation}
  a variant of the adaptive Lasso with text-informed variable-specific
  penalties, nesting the LLM-Lasso approach of
  incorporating domain knowledge from large language models into
  penalized regression.
\end{enumerate}
\end{proposition}

\begin{proof}
\emph{The proof is given in Appendix~\ref{app:proofs}.}
\end{proof}

\begin{remark}[Weak dominance]
\label{rem:weak_dominance}
Since standard KRR corresponds to $\mu_j = 0$ for all $j$,
and the cross-validated theory-informed KRR is free to select
$\hat{\mu}_j = 0$ if theory $j$ is unhelpful, the theory-informed
estimator weakly dominates standard KRR in out-of-sample
prediction loss---up to the granularity of the hyperparameter search.
The same logic applies to each nested special case: the larger
hyperparameter space can only improve (or match) the best
achievable out-of-sample performance.
\end{remark}

%======================================================================
\subsection{Cross-Validation for Hyperparameter Selection}
\label{sub:cv}
%======================================================================

The hyperparameter vector
$\boldsymbol{\lambda} = (\lambda_{\mathrm{stat}}, \mu_1, \ldots, \mu_J)$
is $(J+1)$-dimensional.
We introduce grouped penalties and temporal cross-validation
to make the selection tractable.

\paragraph{Grouped penalties.}
We assign each model $j$ to one of $G$ groups based on its
economic type.  Let $g: \{1, \ldots, J\} \to \{1, \ldots, G\}$
denote the group assignment, and impose $\mu_j = \mu_{g(j)}$
for all $j$.  In our empirical application, $G = 8$, corresponding to
the four penalty types (A--D) each potentially with two variants
(e.g., consumption-based vs.\ production-based SDFs within Type~A).
This reduces the hyperparameter space from $(J+1)$ to $(G+1) = 9$ dimensions.

\paragraph{Temporal cross-validation.}
Financial panel data exhibit strong serial dependence, so we use
leave-one-year-out (LOYO) cross-validation to select $\boldsymbol{\lambda}$.
Partition the sample into $T_{\mathrm{yr}}$ calendar years
$\mathcal{T}_1, \ldots, \mathcal{T}_{T_{\mathrm{yr}}}$.
For fold $\ell$, let
$\mathcal{S}_{-\ell} = \{(i,t) : t \notin \mathcal{T}_\ell\}$ be the
training set and $\mathcal{S}_\ell = \{(i,t) : t \in \mathcal{T}_\ell\}$
the validation set.
The cross-validated objective is
\begin{equation}
\label{eq:cv}
\mathrm{CV}(\boldsymbol{\lambda}) \;=\;
\frac{1}{T_{\mathrm{yr}}} \sum_{\ell=1}^{T_{\mathrm{yr}}}
\frac{1}{|\mathcal{S}_\ell|}
\sum_{s \in \mathcal{S}_\ell}
\bigl(R_s - \hat{f}_{-\ell}(\mathbf{X}_s; \boldsymbol{\lambda})\bigr)^2,
\end{equation}
where $\hat{f}_{-\ell}(\cdot; \boldsymbol{\lambda})$ is the theory-informed
KRR estimator trained on $\mathcal{S}_{-\ell}$ with hyperparameters
$\boldsymbol{\lambda}$.

\paragraph{Search strategy.}
We minimize $\mathrm{CV}(\boldsymbol{\lambda})$ by random search over
a 9-dimensional log-uniform grid:
$\log_{10}\lambda_{\mathrm{stat}} \in [-6, 2]$ and
$\log_{10}\mu_g \in [-6, 2]$ for $g = 1, \ldots, G$.
We draw $B = 500$ random configurations and evaluate each by solving
the finite-dimensional problem~\eqref{eq:finite_dim} $T_{\mathrm{yr}}$
times (once per fold).
The selected hyperparameters are
$\hat{\boldsymbol{\lambda}} = \arg\min_{\boldsymbol{\lambda}}
\mathrm{CV}(\boldsymbol{\lambda})$.

\begin{remark}[Computational considerations]
\label{rem:cv_computation}
Each fold requires solving a linear system of size $|\mathcal{S}_{-\ell}|$,
which is $O(((T_{\mathrm{yr}}-1) \cdot N / T_{\mathrm{yr}})^3)$ per fold.
With $T_{\mathrm{yr}} \approx 30$--$60$ years and $N \approx 3{,}000$--$10{,}000$
assets, this is demanding but feasible with precomputed kernel
eigendecompositions and parallelization across folds and random draws.
The total cost scales as $O(B \cdot T_{\mathrm{yr}} \cdot n_{\mathrm{fold}}^3)$
where $n_{\mathrm{fold}} = |\mathcal{S}_{-\ell}|$.
\end{remark}

%======================================================================
\subsection{Misspecification Robustness}
\label{sub:misspecification}
%======================================================================

A central concern with theory-informed estimation is that imposing
restrictions from a misspecified model may worsen predictions.
The following result shows that cross-validation provides an automatic
safeguard.

\begin{proposition}[Misspecification robustness]
\label{prop:misspecification}
Suppose model $j$ is misspecified in the sense that the true
conditional expected return function $f^*$ satisfies
$C_j(f^*) > 0$.
Under the following regularity conditions:
\begin{enumerate}[label=\textnormal{(R\arabic*)}]
\item \label{cond:consistency}
  The RKHS $\mathcal{H}$ is rich enough that
  $\inf_{f \in \mathcal{H}} \|f - f^*\|_{L^2} = 0$;
\item \label{cond:cv_oracle}
  The cross-validation procedure~\eqref{eq:cv} is asymptotically
  oracle-efficient in the sense that
  $\mathrm{CV}(\hat{\boldsymbol{\lambda}}) - \inf_{\boldsymbol{\lambda}}
  \mathrm{CV}(\boldsymbol{\lambda}) = o_p(1)$;
\item \label{cond:bounded}
  The penalty $C_j$ is continuous in $\|\cdot\|_{L^2}$ and the
  parameter space for $\mu_j$ includes zero;
\end{enumerate}
the cross-validated multiplier satisfies
$\hat{\mu}_j \xrightarrow{p} 0$ as $n \to \infty$.
\end{proposition}

\begin{proof}
\emph{The proof is given in Appendix~\ref{app:proofs}.}
\end{proof}

\begin{remark}[Practical implication]
\label{rem:misspec_practical}
Proposition~\ref{prop:misspecification} provides a theoretical guarantee
that misspecified theories ``turn themselves off'' through cross-validation.
In finite samples, a moderately misspecified theory may still receive
positive $\hat{\mu}_j$ if it captures useful structure despite not
being exactly correct.
The estimator thus exploits partial truths---a theory need not be
exactly right to be useful for prediction.
\end{remark}

\begin{remark}[Comparison to Bayesian model averaging]
\label{rem:bma}
The misspecification robustness of theory-informed KRR is analogous to
Bayesian model averaging (BMA) in spirit: both frameworks downweight
models that fit poorly out of sample.
However, the mechanism differs.
BMA assigns posterior probabilities to discrete model specifications,
whereas theory-informed KRR treats each model's restrictions as a
continuous penalty with cross-validated strength.
The latter avoids the combinatorial explosion of the BMA model space
and allows partial incorporation of each theory's restrictions.
\end{remark}

%======================================================================
\subsection{Inference on Model Value}
\label{sub:inference}
%======================================================================

Beyond prediction, one may wish to test whether a particular economic
model contributes statistically significant information beyond agnostic
regularization.
This amounts to testing
\begin{equation}
\label{eq:test_mu}
H_0: \mu_j^* = 0
\qquad \text{vs.} \qquad
H_1: \mu_j^* > 0,
\end{equation}
where $\mu_j^*$ is the population-optimal multiplier for model $j$.
Rejection of $H_0$ implies that model $j$ provides marginal predictive
value beyond what the kernel and other models already capture.

%----------------------------------------------------------------------
\subsubsection{Connection to Uniform Inference for KRR}
\label{sub:uniform_inference}
%----------------------------------------------------------------------

Recent work by \citet*{SinghVijaykumar2023} establishes uniform
confidence bands for kernel ridge regression estimators of conditional
expectations.
Their framework provides pointwise and uniform inference for
$f^*(\mathbf{x})$ across $\mathbf{x} \in \mathcal{X}$ using
a bias-corrected estimator and variance estimation based on the
kernel's spectral properties.

We outline how this can be extended to test~\eqref{eq:test_mu}.

\begin{proposition}[Testing the marginal value of economic models]
\label{prop:test_mu}
Suppose the conditions of \citet*{SinghVijaykumar2023} hold for
the RKHS $\mathcal{H}$, augmented with the following:
\begin{enumerate}[label=\textnormal{(T\arabic*)}]
\item \label{cond:test_smooth}
  The mapping $\mu_j \mapsto \hat{f}(\cdot\,; \mu_j)$ is differentiable
  in $L^2$, with derivative $\partial_{\mu_j} \hat{f}$ satisfying
  $\|\partial_{\mu_j} \hat{f}\|_{L^2} = O(1)$;
\item \label{cond:test_rate}
  The estimation error satisfies
  $\|\hat{f}(\cdot\,; \hat{\boldsymbol{\lambda}}) - f^*\|_\infty = O_p(r_n)$
  for a rate $r_n \to 0$.
\end{enumerate}
Define the test statistic
\begin{equation}
\label{eq:T_stat}
T_j \;=\;
n \cdot \bigl[\mathrm{CV}(\hat{\boldsymbol{\lambda}}_{-j})
- \mathrm{CV}(\hat{\boldsymbol{\lambda}})\bigr],
\end{equation}
where $\hat{\boldsymbol{\lambda}}_{-j}$ denotes the cross-validated
hyperparameters with $\mu_j$ constrained to zero.
Under $H_0\!: \mu_j^* = 0$, $T_j = o_p(1)$.
Under $H_1\!: \mu_j^* > 0$, $T_j \to \infty$ in probability.
\end{proposition}

\begin{proof}
\emph{The proof is given in Appendix~\ref{app:proofs}.}
\end{proof}

\begin{remark}[Permutation inference]
\label{rem:permutation}
A full distributional theory for $T_j$ under $H_0$ (critical values,
size control, power analysis) requires characterizing the joint
distribution of the cross-validated hyperparameters, which we defer
to future work.
In practice, we implement permutation-based inference:
randomly reassign the structural penalty $C_j$ to pseudo-models
(scrambling the economic content while preserving the penalty's
functional form) and compare $T_j$ to the permutation distribution.
This yields valid $p$-values under exchangeability of the
pseudo-model assignments.
\end{remark}

\begin{remark}[Interpretation]
\label{rem:test_interpretation}
The test~\eqref{eq:test_mu} asks a fundamentally different question
from testing whether a model's pricing errors are zero.
A model may have large pricing errors ($C_j(f^*) > 0$) yet still
receive $\mu_j^* > 0$ if its restrictions are correlated with the
true $f^*$ in a way that improves finite-sample estimation.
Conversely, a correctly specified model ($C_j(f^*) = 0$) may have
$\mu_j^* = 0$ if the RKHS and data are rich enough that the model's
restrictions are redundant.
The test thus measures the \emph{marginal predictive value} of each
theory in the presence of all others.
\end{remark}

\begin{remark}[Relation to model confidence sets]
\label{rem:mcs}
The pairwise comparison underlying~\eqref{eq:T_stat} connects to
the model confidence set (MCS) framework of
\citet*{HansenLundeNason2011}.
One could construct an MCS of theories whose marginal contributions
are statistically indistinguishable:
$\widehat{\mathrm{MCS}}_\alpha = \{j : T_j \leq c_\alpha\}$,
where $c_\alpha$ is calibrated by the permutation distribution.
This provides a principled way to report which theories are
``useful'' for asset-pricing prediction at a given confidence level.
\end{remark}

%======================================================================
\subsection{Summary of the Framework}
\label{sub:summary_framework}
%======================================================================

Table~\ref{tab:framework_summary} collects the main components of the
theory-informed KRR estimator.

\begin{table}[htbp]
\centering
\caption{Summary of the theory-informed KRR framework}
\label{tab:framework_summary}
\small
\begin{tabular}{@{}lll@{}}
\toprule
\textbf{Component} & \textbf{Role} & \textbf{Reference} \\
\midrule
Prediction loss & Fit to observed returns & Eq.~\eqref{eq:tikrr} \\
RKHS penalty & Agnostic statistical regularization & Eq.~\eqref{eq:tikrr} \\
Type~A penalty & Euler equation restrictions (SDF models) & Def.~\ref{def:euler_penalty} \\
Type~B penalty & Production/investment FOC restrictions & Def.~\ref{def:prod_penalty} \\
Type~C penalty & Factor structure restrictions & Def.~\ref{def:factor_penalty} \\
Type~D penalty & Demand-system equilibrium restrictions & Def.~\ref{def:demand_penalty} \\
Representer theorem & Finite-dimensional reduction & Prop.~\ref{prop:representer} \\
First-order conditions & Closed-form for quadratic penalties & Prop.~\ref{prop:foc} \\
Nesting & Recovers KRR, GMM, IPCA, Lasso variants & Prop.~\ref{prop:nesting} \\
CV selection & Grouped penalties, temporal LOYO & Eq.~\eqref{eq:cv} \\
Misspecification & Misspecified theories auto-downweight & Prop.~\ref{prop:misspecification} \\
Inference & Test marginal value of each theory & Prop.~\ref{prop:test_mu} \\
\bottomrule
\end{tabular}
\end{table}

%!TEX root = ../main.tex
%----------------------------------------------------------------------
%  Section 3: Structural Restrictions Catalog
%  Paper:   Theory-Informed Kernel Ridge Regression for Asset Pricing
%  Authors: Amedeo Andriollo and Juan F. Imbet
%----------------------------------------------------------------------

\section{Structural Restrictions Catalog}\label{sec:restrictions}

Each restriction in our framework originates from a specific economic model and imposes a distinct functional constraint on the mapping from observable characteristics to expected returns.
A consumption-based Euler equation with power utility implies a different penalty $C_j(f)$ than one with Epstein--Zin preferences, even though both restrict the same SDF moment condition, because the functional form of the stochastic discount factor $M_{j,t+1}$ differs.
Similarly, a production-based investment CAPM penalty differs from a $q$-theory penalty not merely in its parameter values but in which investment variables enter the mapping $g_j$ and how they interact.
The penalties are not interchangeable: each encodes a qualitatively different economic hypothesis about why expected returns vary across assets and over time.

We organize the fifty-six restrictions into eight families, grouped by economic mechanism: consumption-based Euler equations, production-based and investment restrictions, intermediary and institutional restrictions, information and learning restrictions, demand-system and market microstructure restrictions, variance and volatility restrictions, behavioral and sentiment restrictions, and macro-finance and term structure restrictions.
For each restriction, we specify three elements: the source model (the economic theory that implies it), the mathematical form of the restriction (the SDF, first-order condition, or equilibrium condition), and the penalty functional $C_j(f)$ through which it enters the Theory-KRR objective~\eqref{eq:tikrr} from Section~\ref{sec:framework}.
Because every penalty $C_j(f)$ depends on $f$ only through its evaluations at observed data points $\{f(\mathbf{X}_s)\}_{s=1}^n$, the representer theorem (Proposition~\ref{prop:representer}) applies to the full collection, and the estimator remains computationally tractable regardless of how many restrictions are active.


%======================================================================
\subsection{Consumption-Based Euler Equations}\label{sub:euler_restrictions}
%======================================================================

Consumption-based models derive asset-pricing restrictions from the representative agent's intertemporal optimality condition.
The Euler equation $\E[M_{t+1} R_{i,t+1}^{\mathrm{gross}} \mid \Fcal_t] = 1$ must hold for every asset $i$ and date $t$, where the SDF $M_{t+1}$ is determined by the agent's preferences and the consumption process.
Different preference specifications---power utility, recursive utility, habit formation, long-run risk---yield different SDFs and hence different penalties.
All thirteen restrictions in this family are Type~A (Euler equation) penalties as defined in Section~\ref{sub:euler}:
\begin{equation}\label{eq:euler_family}
C_j^{\mathrm{Euler}}(f) \;=\;
\sum_{t=1}^{T}
\left(
\frac{1}{N}\sum_{i=1}^{N}
M_j\bigl(\hat{\theta}_j, \mathbf{X}_t, \mathbf{X}_{t+1}\bigr)
\cdot
\bigl(f(\mathbf{X}_{i,t}) + R_t^f\bigr) - 1
\right)^2,
\end{equation}
where $\hat{\theta}_j$ denotes pre-estimated preference parameters obtained by GMM.
What distinguishes the thirteen restrictions is solely the form of $M_j$.

Table~\ref{tab:euler} lists the thirteen consumption-based restrictions. The thirteen consumption-based restrictions span several distinct economic channels.
Restrictions~1--5 vary preferences while maintaining i.i.d.\ or simple Markov consumption dynamics; what differs is how marginal utility depends on consumption growth and, in the habit models, on the history of consumption.
Power utility (Restriction~1) serves as the baseline, and each subsequent preference specification adds a mechanism---intertemporal substitution separated from risk aversion (Restriction~2), slow-moving external habit (Restriction~3), or direct dependence on lagged consumption (Restrictions~4--5).
These five SDFs share the same data requirements (aggregate consumption growth) but impose different functional relationships between $M_{t+1}$ and $\Delta c_{t+1}$.

Restrictions~6--9 modify the consumption \emph{process} rather than preferences alone.
Long-run risk (Restriction~6) introduces a persistent expected-growth component $x_t$ and stochastic volatility $\sigma_t$, both of which enter the SDF through Epstein--Zin recursion; the SDF depends on the news to $x_t$ and $\sigma_t^2$, not merely on realized consumption growth.
Adding jumps (Restriction~7) generates occasional large movements in volatility that are absent under Gaussian dynamics.
Rare disaster models (Restrictions~8--9) replace continuous volatility variation with low-probability, high-impact consumption contractions.
The economic content differs sharply: long-run risk explains the equity premium through intertemporal hedging demand, while disaster models explain it through the pricing of tail events.

Restrictions~10--13 extend the framework along further dimensions.
Precautionary savings (Restriction~10) introduces higher-order risk terms via the prudence coefficient, generating SDF dependence on conditional consumption variance.
Durable consumption (Restriction~11) adds a second consumption good with its own intratemporal elasticity.
Heterogeneous agents (Restriction~12) replaces the representative consumer with a continuum of agents whose idiosyncratic consumption risk enters the aggregate SDF through the cross-sectional variance of consumption growth.
The nonparametric Euler restriction (Restriction~13) abandons parametric preferences entirely, requiring only that some monotone decreasing function of consumption growth prices assets correctly---it serves as a benchmark that detects whether \emph{any} consumption-based SDF is empirically useful, without committing to a specific functional form.


%======================================================================
\subsection{Production-Based and Investment Restrictions}\label{sub:prod_restrictions}
%======================================================================

Production-based asset-pricing models derive expected returns from firms' intertemporal investment decisions.
The key insight is that firms equate the marginal cost of investment today to the expected discounted marginal benefit, and this first-order condition links expected stock returns to observable investment and profitability variables.
All ten restrictions in this family are Type~B (production FOC) penalties as defined in Section~\ref{sub:production}:
\begin{equation}\label{eq:prod_family}
C_j^{\mathrm{prod}}(f) \;=\;
\frac{1}{n}\sum_{s=1}^{n}
\bigl(f(\mathbf{X}_s) - g_j(\mathbf{Z}_s^{(j)})\bigr)^2,
\end{equation}
where $g_j$ maps firm-level investment characteristics $\mathbf{Z}_{i,t}^{(j)}$ to expected returns.
The functional form of $g_j$ and the variables that enter $\mathbf{Z}_{i,t}^{(j)}$ vary across restrictions.

Table~\ref{tab:production} lists the ten production-based restrictions. The production-based restrictions divide into three groups.
Restrictions~14--18 concern the technology of physical capital adjustment.
The investment CAPM (Restriction~14) and $q$-theory (Restriction~15) provide the broadest frameworks: expected returns are determined by firms' optimal investment behavior, with investment-to-capital ratio and profitability as the key state variables.
Restrictions~16--18 refine the adjustment cost specification.
Convex adjustment costs (Restriction~16) yield a smooth, monotone relationship between $I/K$ and expected returns.
Nonconvex adjustment costs (Restriction~17) and irreversibility (Restriction~18) break this smoothness: large adjustments incur fixed costs, and disinvestment is more expensive than investment.
These restrictions generate different functional forms for $g_j$---Restriction~16 predicts a linear or quadratic relationship in $I/K$, while Restrictions~17--18 predict threshold effects and asymmetries.

Restrictions~19--21 address financial frictions.
These models predict that expected returns reflect not only the marginal product of capital but also the shadow cost of financing constraints.
The three restrictions differ in which friction binds: collateral constraints (Restriction~19) link the premium to the ratio of capital to debt, equity issuance costs (Restriction~20) link it to leverage and internal cash flow, and costly external finance (Restriction~21) links it to the dividend payout and leverage.
All three predict that financially constrained firms earn higher expected returns, but the mapping from observables to the constraint's shadow value differs across models.

Restrictions~22--23 expand the set of productive inputs beyond physical capital.
Labor adjustment costs (Restriction~22) introduce the hiring rate as an additional state variable in the investment FOC, generating a two-dimensional mapping $(I/K, \mathrm{HN}) \mapsto \E[R]$.
Intangible capital (Restriction~23) adds R\&D spending and selling, general, and administrative expenses as measures of intangible investment, recognizing that knowledge capital and organizational capital are productive assets with their own adjustment costs.


%======================================================================
\subsection{Intermediary and Institutional Restrictions}\label{sub:intermediary_restrictions}
%======================================================================

Intermediary asset-pricing models posit that marginal investors are not households but financial intermediaries---broker-dealers, banks, hedge funds---whose pricing kernels reflect institutional constraints on leverage, capital, and funding.
When intermediaries are the marginal pricers, expected returns depend on intermediary-specific state variables: capital ratios, leverage, funding spreads, and balance-sheet capacity.
A linear factor model with factor $F_{j,t}$ implies the SDF $M_{j,t+1} = a_j + b_j F_{j,t+1}$, which is a special case of a linear SDF.
All eight restrictions in this family are therefore Type~A (Euler equation) penalties as defined in Section~\ref{sub:euler}:
\begin{equation}\label{eq:intermediary_family}
C_j^{\mathrm{Euler}}(f) \;=\;
\sum_{t=1}^{T}
\left(
\frac{1}{N}\sum_{i=1}^{N}
\bigl(\hat{a}_j + \hat{b}_j\, F_{j,t+1}\bigr)
\cdot
\bigl(f(\mathbf{X}_{i,t}) + R_t^f\bigr) - 1
\right)^2,
\end{equation}
where $F_{j,t}$ is the intermediary factor and $\hat{a}_j$, $\hat{b}_j$ are pre-estimated SDF parameters obtained by GMM.
This formulation requires only the factor time series---no pre-estimated betas are needed.

Table~\ref{tab:intermediary} lists the eight intermediary restrictions. The intermediary restrictions separate into models of intermediary health (Restrictions~24--25, 27) and models of market frictions that constrain intermediation (Restrictions~26, 28--31).
The intermediary capital ratio (Restriction~24) and leverage (Restriction~25) factors capture the state of broker-dealer balance sheets: when capital is depleted or leverage is high, the marginal value of wealth to the intermediary rises, increasing risk prices.
The dealer balance sheet factor (Restriction~27) constructs an equity capital ratio specifically for primary dealers, providing a more direct measure of the capacity constraint.
These three restrictions use different empirical proxies for the same conceptual channel---intermediary distress---and their relative $\mu_j$ values reveal which proxy captures the pricing-relevant variation.

Restrictions~26 and 28--31 focus on frictions in the trading environment.
Funding liquidity (Restriction~26) links expected returns to the cost and availability of short-term funding: when funding dries up, margins increase and leveraged positions are unwound, creating fire-sale dynamics.
Limits to arbitrage (Restriction~28) and slow-moving capital (Restriction~29) introduce a role for mispricing persistence---arbitrageurs face constraints (capital, risk limits, organizational frictions) that prevent immediate price correction.
The margin CAPM (Restriction~30) derives a two-beta model in which both market exposure and margin exposure are priced, while the leverage-constraint model (Restriction~31) predicts a compressed security market line because leverage-constrained investors overweight high-beta assets.
The key distinction: Restrictions~24--27 model \emph{who} the marginal pricer is, while Restrictions~28--31 model \emph{what constrains} the marginal pricer.


%======================================================================
\subsection{Information and Learning Restrictions}\label{sub:info_restrictions}
%======================================================================

Information and learning models generate expected-return variation from agents' incomplete knowledge about economic fundamentals.
Unlike consumption-based models, where the SDF is pinned down by preferences and the consumption process, these models derive pricing implications from the structure of beliefs---how agents learn, what they are uncertain about, and how they process information.
These restrictions generate Type~A penalties: Restrictions~32--35 specify a complete SDF from the learning model, while Restrictions~36--37 imply a linear SDF driven by a belief-related state variable (disagreement or information flow), which reduces to the Euler equation form $M_{j,t+1} = a_j + b_j F_{j,t+1}$.

Table~\ref{tab:information} lists the six information restrictions. The six information restrictions differ in \emph{what} agents are uncertain about and \emph{how} they process that uncertainty.
Bayesian learning (Restriction~32) generates time-varying risk premia because the posterior variance of growth declines as data accumulate, changing the conditional distribution of the SDF.
Regime learning (Restriction~33) produces a distinct mechanism: agents infer a latent state (expansion vs.\ recession), and uncertainty about the current regime amplifies return volatility and risk premia---the SDF depends on the belief $\pi_t$ rather than on a parameter estimate.
Parameter uncertainty (Restriction~34) prices the risk that the true data-generating process is unknown, generating a premium for model risk that shrinks only slowly with sample size.

Ambiguity aversion (Restriction~35) departs from Bayesian updating altogether: the agent entertains a set of plausible models $\mathcal{P}$ and evaluates each asset under the worst-case model.
The SDF is the likelihood ratio between the worst-case and reference models, generating pricing implications distinct from those of any single Bayesian model.
Disagreement (Restriction~36) and rational inattention (Restriction~37) shift focus from a single agent's beliefs to the cross-section of beliefs.
Disagreement models price the dispersion of beliefs as a risk factor (optimists hold long positions, pessimists short, and their trading generates a risk premium).
Rational inattention models generate endogenous information asymmetry: agents with limited processing capacity allocate attention optimally, and the resulting heterogeneity in information sets produces cross-sectional variation in expected returns.


%======================================================================
\subsection{Demand-System and Market Microstructure Restrictions}\label{sub:demand_restrictions}
%======================================================================

Demand-system models derive asset prices from the aggregation of heterogeneous investors' portfolio demands.
The equilibrium restriction is market clearing: aggregate demand for each asset must equal its supply.
This yields Type~C (demand-system equilibrium) penalties as defined in Section~\ref{sub:demand}:
\begin{equation}\label{eq:demand_family}
C_j^{\mathrm{demand}}(f) \;=\;
\sum_{t=1}^{T}
\left\|
\sum_{i} w_i^{(j)}\, D_i^{(j)}\bigl(\hat{\mathbf{f}}_t\bigr) - \mathbf{S}_t
\right\|^2,
\end{equation}
where $\hat{\mathbf{f}}_t = (f(\mathbf{X}_{1,t}), \ldots, f(\mathbf{X}_{N,t}))^\top$ and $\mathbf{S}_t$ denotes market-capitalization shares.
Restrictions~41--43 involve microstructure frictions and take modified forms as described below.

Table~\ref{tab:demand} lists the six demand-system and microstructure restrictions. The demand-system restrictions (38--40) share the market-clearing equilibrium condition but differ in the structure of investors' demand functions.
Restriction~38 models demand as a function of observable characteristics and a latent demand shock, estimating the demand system directly from portfolio holdings data.
Restriction~39 derives demand from investors' optimality conditions with characteristic-based demand elasticities, providing a micro-founded demand function.
Restriction~40 focuses on the \emph{aggregate} demand elasticity: when total market demand is inelastic, flow-driven trading has outsized price impact, and expected returns reflect the accommodation of demand shocks rather than fundamental risk.

Restrictions~41--43 address frictions that prevent instantaneous and costless trading.
Search frictions (Restriction~41) apply primarily to over-the-counter markets, where finding a counterparty takes time; the expected-return penalty includes the search intensity and the buyer's bargaining power.
Fire sale externalities (Restriction~42) generate a Type~A Euler penalty with $M_{42,t+1} = \hat{a}_{42} + \hat{b}_{42}\,\mathrm{FSP}_{t+1}$: when multiple leveraged institutions simultaneously liquidate positions, the fire-sale pressure factor enters the SDF.
Short-sale constraints (Restriction~43) restrict the adjustment mechanism on the negative side: when pessimists cannot short, asset prices reflect only optimists' valuations, and the degree of overpricing depends on the divergence of opinions and the binding intensity of the constraint.


%======================================================================
\subsection{Variance and Volatility Restrictions}\label{sub:option_restrictions}
%======================================================================

Variance and volatility models link expected returns to measures of conditional variance and the compensation investors demand for bearing variance risk.
These two restrictions use publicly available data---the CBOE Volatility Index (VIX) as a measure of implied variance and daily CRSP returns to compute realized variance---rather than individual equity option data, making them applicable to the broad cross-section over the post-1990 period.
Both restrictions generate Type~B penalties that map volatility-related state variables to expected returns.

Table~\ref{tab:options} lists the two variance and volatility restrictions. The two variance and volatility restrictions capture complementary aspects of volatility pricing.
The variance risk premium (Restriction~44) uses the spread between implied variance (from VIX) and expected realized variance (computed from daily CRSP returns within each month) as a predictor of aggregate returns; the penalty penalizes discrepancies between $f(\mathbf{X}_t)$ and the VRP-implied expected return.
\citet{bollerslevtauchenzhou2009} show that VRP is a strong predictor of market returns at horizons from one to six months, and the restriction encodes this predictive relationship.
Stochastic volatility (Restriction~45) uses the conditional variance $V_t$ from a Heston-type model, where VIX serves as a proxy for the latent conditional variance and the variance process carries its own risk premium distinct from the equity premium.
Together, these restrictions ensure that the estimator respects the well-documented relationship between variance dynamics and expected returns without requiring granular option-level data.


%======================================================================
\subsection{Behavioral and Sentiment Restrictions}\label{sub:behavioral_restrictions}
%======================================================================

Behavioral models derive expected-return patterns from departures from rational expectations or expected utility.
These restrictions do not derive an SDF from a representative agent's Euler equation; instead, they predict cross-sectional or time-series return patterns as functions of behavioral variables---past returns, capital gains, sentiment indices.
Most generate Type~B penalties, where $g_j$ maps behavioral state variables to expected returns.

Table~\ref{tab:behavioral} lists the six behavioral restrictions. The behavioral restrictions exploit distinct psychological mechanisms.
Prospect theory (Restriction~46) modifies the SDF with a kink at the reference point: losses are weighted more heavily than gains by a factor $\lambda > 1$, generating a first-order risk premium for assets whose returns covary with movements around the reference point.
This is the only behavioral restriction that directly specifies an SDF, placing it closest to the consumption-based framework.

Restrictions~47--49 generate cross-sectional return predictability from belief formation biases.
Overconfidence (Restriction~47) produces momentum at short horizons (agents overreact to private signals, pushing prices beyond fundamentals) and reversal at long horizons (prices revert as public information arrives).
The disposition effect (Restriction~48) predicts return continuation among stocks with large unrealized gains and return reversal among stocks with large unrealized losses, because investors' reluctance to realize losses impedes price adjustment.
Extrapolative expectations (Restriction~49) model agents who weight recent returns too heavily when forming expectations, generating a specific functional form for the relationship between past returns and expected future returns.
These three restrictions share momentum-related variables as inputs but predict different functional relationships between past returns and expected future returns.

Investor sentiment (Restriction~50) operates at the aggregate level: the Baker--Wurgler sentiment index predicts that, conditional on high sentiment, speculative stocks (small, unprofitable, high-volatility) earn low subsequent returns, while safe stocks are largely unaffected.
This generates a Type~A Euler equation penalty where $M_{50,t+1} = \hat{a}_{50} + \hat{b}_{50}\,\mathrm{Sent}_{t+1}$ is a linear SDF in the sentiment factor.
Salience theory (Restriction~51) predicts mispricing at the individual asset level: stocks with salient (extreme) upside payoffs are overpriced because agents overweight the probability of salient outcomes, generating a negative relationship between salience and expected returns.


%======================================================================
\subsection{Macro-Finance and Term Structure Restrictions}\label{sub:macro_restrictions}
%======================================================================

Macro-finance models link expected returns on equities and bonds to macroeconomic state variables: the yield curve, inflation expectations, monetary policy, and consumption-wealth ratios.
These restrictions enter as Type~A or Type~B penalties depending on whether the model implies a linear SDF driven by a macro factor (Type~A) or a direct functional mapping from macro variables to expected returns (Type~B).

Table~\ref{tab:macro} lists the five macro-finance restrictions. The macro-finance restrictions link equity expected returns to macroeconomic and fixed-income state variables through five distinct channels.
The affine term structure restriction (Restriction~52) imposes no-arbitrage across maturities, using the yield curve state vector $\mathbf{y}_t$ (level, slope, curvature, and potentially macro factors) to constrain the time-varying equity risk premium.
This restriction is structurally different from the others because it derives from the bond market and imposes cross-market no-arbitrage.

Restrictions~53--54 transmit monetary and inflation shocks to equity prices.
The inflation risk premium (Restriction~53) uses the spread between nominal and real yields to extract the compensation investors demand for bearing inflation uncertainty; this premium predicts equity returns because inflation uncertainty covaries with real economic risk.
The monetary policy restriction (Restriction~54) identifies the surprise component of federal funds rate decisions (using fed funds futures) and maps it to equity returns, exploiting the direct channel from policy rates to discount rates.

Restrictions~55 and~56 use slow-moving macroeconomic aggregates.
The consumption-wealth ratio $cay$ (Restriction~55) exploits cointegration among consumption, asset wealth, and labor income: when consumption is high relative to wealth, expected returns are high because the ratio mean-reverts through subsequent returns.
The excess bond premium (Restriction~56) isolates the component of credit spreads that reflects risk appetite and financial intermediary health, net of default risk; it predicts equity returns because deteriorating credit conditions signal tightening financial constraints economy-wide.


%======================================================================
\subsection{Summary and Discussion}\label{sub:restriction_summary}
%======================================================================

The catalog comprises fifty-six restrictions organized into eight families.
Each restriction imposes a distinct functional form on the mapping from observable characteristics $\mathbf{X}_{i,t}$ to expected returns $f(\mathbf{X}_{i,t})$, and each translates into a specific penalty functional $C_j(f)$ within the Theory-KRR objective~\eqref{eq:tikrr}.
The penalties are heterogeneous in structure: Type~A penalties enforce Euler equation moment conditions and require specifying the SDF---either from a structural model or from a linear factor representation $M_{t+1} = a + b F_{t+1}$ (Restrictions~1--13, 24--37, 42, 46, 50, 53, 56); Type~B penalties equate expected returns to functions of investment, behavioral, or macro variables through first-order conditions (Restrictions~14--23, 41, 43, 44--45, 47--49, 51--52, 54--55); and Type~C penalties enforce market-clearing equilibrium conditions from demand-system models (Restrictions~38--40).
Despite this heterogeneity, the framework in Section~\ref{sec:framework} accommodates all fifty-six penalties within a single objective: each $C_j(f)$ depends on $f$ only through its evaluations at observed data points, the representer theorem applies, and cross-validation selects the vector $\bm{\mu}$ that balances predictive fit against theoretical guidance.

A property of the catalog that makes the $\mu_j$ ranking meaningful is that many restrictions share observable input variables but impose different functional relationships.
Consumption growth enters Restrictions~1--13 through thirteen different SDFs.
Investment-to-capital ratio enters Restrictions~14--18 through five different adjustment cost specifications.
Past returns enter Restrictions~47--49 through three different behavioral mechanisms.
When two restrictions share inputs but one receives a higher $\mu_j$ in cross-validation, the data have identified which functional relationship between those inputs and expected returns is more empirically useful.
This is the sense in which the cross-validated multipliers $\bm{\mu}$ produce a ranking of economic theories: they rank not just which state variables matter, but which \emph{economic mechanisms linking state variables to expected returns} are supported by the data.

%!TEX root = ../main.tex
%----------------------------------------------------------------------
%  Section 4: Estimation and Implementation
%  Paper:   Theory-Informed Kernel Ridge Regression for Asset Pricing
%  Authors: Amedeo Andriollo and Juan F. Imbet
%----------------------------------------------------------------------

\section{Estimation and Implementation}\label{sec:estimation}

This section addresses the practical aspects of taking Theory-KRR from the mathematical framework of Sections~\ref{sec:framework}--\ref{sec:restrictions} to an implementable estimator.
We describe the kernel choice and bandwidth calibration (Section~\ref{sub:kernel_choice}), the pre-estimation of structural parameters that enter the penalty functionals (Section~\ref{sub:pre_estimation}), the cross-validation strategy for selecting the high-dimensional vector of penalty multipliers (Section~\ref{sub:cv_strategy}), computational details (Section~\ref{sub:computation}), and the rolling-window procedure for out-of-sample evaluation (Section~\ref{sub:rolling}).

%======================================================================
\subsection{Kernel Choice and Bandwidth Selection}\label{sub:kernel_choice}
%======================================================================

The kernel $k: \R^p \times \R^p \to \R$ determines the RKHS $\Hcal$ in which the estimator searches for the conditional expectation function $f^*$.
Our primary specification is the Gaussian radial basis function (RBF) kernel:
\begin{equation}\label{eq:gaussian_kernel}
k(\mathbf{x}, \mathbf{x}') = \exp\!\left(-\frac{\|\mathbf{x} - \mathbf{x}'\|^2}{2\sigma^2}\right),
\end{equation}
where $\sigma > 0$ is the bandwidth parameter and $\|\cdot\|$ denotes the Euclidean norm in $\R^p$.
The Gaussian kernel is universal: the associated RKHS is dense in the space of continuous functions on any compact subset of $\R^p$, so the approximation error $\inf_{f \in \Hcal}\|f - f^*\|$ can be made arbitrarily small.
This universality is essential for our application, where the true conditional expectation function may exhibit the nonlinearities, interactions, and threshold effects documented in the cross-section of returns.

We set the bandwidth $\sigma$ using the median heuristic:
\begin{equation}\label{eq:median_heuristic}
\sigma = \operatorname{median}\!\left\{\|\mathbf{X}_{i,t} - \mathbf{X}_{j,s}\| : (i,t) \neq (j,s)\right\}.
\end{equation}
The median heuristic places the kernel bandwidth at a scale where roughly half of the pairwise distances produce kernel evaluations above $e^{-1/2} \approx 0.61$ and half below.
This data-driven rule avoids both extremes: a bandwidth that is too small relative to the data's spread produces a nearly diagonal kernel matrix (overfitting), while a bandwidth that is too large produces a nearly constant kernel matrix (underfitting, approaching the global mean).
The median heuristic has strong theoretical support in the kernel methods literature and performs well in high-dimensional settings where manual tuning is infeasible.

Two limiting cases clarify the role of the kernel.
A linear kernel $k(\mathbf{x}, \mathbf{x}') = \mathbf{x}^\top \mathbf{x}'$ restricts $\Hcal$ to linear functions $f(\mathbf{x}) = \bm{\beta}^\top \mathbf{x}$, and Theory-KRR reduces to a penalized linear regression with structural penalties on coefficients.
This nesting (Proposition~\ref{prop:nesting}(iii)) allows us to isolate the contribution of nonlinearity from the contribution of theory-informed penalties.
At the other extreme, the Gaussian kernel with bandwidth $\sigma \to 0$ produces a kernel matrix approaching the identity, so $\hat{f}$ interpolates the data subject only to the structural penalties.
The joint configuration of $\sigma$, $\lambda_{\mathrm{stat}}$, and $\bm{\mu}$ determines the estimator's regularization, and cross-validation selects the balance that minimizes OOS prediction error.

%======================================================================
\subsection{Estimation of Structural Parameters}\label{sub:pre_estimation}
%======================================================================

Each structural penalty $C_j(f)$ depends on model-specific parameters $\theta_j$ that characterize the underlying economic model.
For consumption-based models, $\theta_j$ includes the coefficient of relative risk aversion $\gamma$, the elasticity of intertemporal substitution $\psi$, the subjective discount factor $\delta$, and the habit persistence parameter $b$.
For production-based models, $\theta_j$ governs the curvature of the adjustment cost function.
For intermediary models, $\theta_j$ parameterizes the SDF loadings on the intermediary factor.
We consider three approaches to estimating $\theta_j$, ordered from most to least principled.

\paragraph{Approach 1: Joint estimation (preferred).}
The most principled approach estimates $f$ and $\{\theta_j\}_{j=1}^J$ simultaneously:
\begin{equation}\label{eq:joint_estimation}
(\hat{f}, \hat{\theta}_1, \ldots, \hat{\theta}_J) \;=\; \argmin_{f \in \mathcal{H}, \, \theta_1, \ldots, \theta_J} \;
\frac{1}{n}\sum_s (R_s - f(\mathbf{X}_s))^2 + \lambda_{\mathrm{stat}}\|f\|_{\mathcal{H}}^2 + \sum_j \mu_j\, C_j(f; \theta_j).
\end{equation}
This optimization is convex in $f$ (given $\theta$) and generally nonconvex in $\theta$, but alternating minimization---updating $f$ given $\theta$ and then $\theta$ given $f$---converges reliably from the pre-estimation starting point.
Joint estimation is the most principled approach because it allows the structural parameters and the nonparametric function to adapt to each other: a consumption-based SDF with risk aversion $\gamma$ that is optimal in isolation may not be optimal when combined with intermediary and production penalties.

\paragraph{Approach 2: Pre-estimation (two-stage).}
For computational simplicity, one can pre-estimate $\hat{\theta}_j$ in a first stage that is separate from the KRR optimization.
For each consumption-based SDF model $j$, we estimate
\begin{equation}\label{eq:gmm_first_stage}
\hat{\theta}_j = \argmin_{\theta_j} \left\|\frac{1}{T}\sum_{t=1}^{T} \mathbf{z}_t \left[M_j(\theta_j, \mathbf{X}_t, \mathbf{X}_{t+1})\, R_{m,t+1}^{\mathrm{gross}} - 1\right]\right\|_{\hat{\mathbf{W}}^{-1}}^2,
\end{equation}
where $\mathbf{z}_t$ is a vector of instruments (lagged consumption growth, the dividend-price ratio, the term spread) and $\hat{\mathbf{W}}$ is an HAC-consistent estimate of the instrument covariance matrix.
For production-based models, we estimate adjustment cost parameters from firm-level investment regressions.
For intermediary models, the SDF loadings $(\hat{a}_j, \hat{b}_j)$ are estimated by GMM on the Euler equation $\E[M_j R^{\mathrm{gross}}] = 1$ using test assets.
The pre-estimated $\hat{\theta}_j$ is then plugged into $C_j(f)$ and treated as fixed during the second-stage KRR optimization.
This approach ignores the interaction between $f$ and $\theta$: the structural parameters estimated from aggregate moments may not be optimal for cross-sectional return prediction.

\paragraph{Approach 3: Calibration.}
For parameters with well-established values, we adopt published calibrations from the original papers.
We set $\delta = 0.99$ (monthly) in all consumption-based models and use $b = 0.87$ from \citet{campbellcochrane1999} for the external habit model.
This is the simplest approach: $C_j(f)$ is fully determined by the data and a known parameter vector, with no estimation of $\theta_j$.
The risk is that published parameter values, calibrated to match aggregate moments in a specific sample, may not be appropriate for cross-sectional return prediction.

Our baseline uses joint estimation (Approach~1), initialized from the pre-estimation values (Approach~2).
Robustness analysis compares all three approaches.
For the sensitivity analysis, we report results over a grid of values for the most influential parameters:
\begin{equation}\label{eq:sensitivity_grid}
\gamma \in \{2, 5, 10, 20\}, \qquad \delta \in \{0.95, 0.97, 0.99\}.
\end{equation}
This grid spans the range of calibrations used in the asset-pricing literature, from low-risk-aversion economies ($\gamma = 2$) to the high-risk-aversion calibrations required by rare disaster models ($\gamma = 20$).

%======================================================================
\subsection{Cross-Validation Strategy}\label{sub:cv_strategy}
%======================================================================

The full hyperparameter vector of Theory-KRR is $(\lambda_{\mathrm{stat}}, \mu_1, \ldots, \mu_J)$, where $J = 56$ in our baseline specification.
A 57-dimensional grid search is computationally infeasible.
We reduce dimensionality by exploiting the economic structure of the restrictions and adopt the three-way temporal split of \citet{gu2020empirical} (Section~1.1) for hyperparameter selection.

We group the 56 restrictions into $G = 8$ families corresponding to the eight economic model classes cataloged in Section~\ref{sec:restrictions}: consumption-based Euler equations, production-based restrictions, intermediary restrictions, information and learning restrictions, demand-system restrictions, variance and volatility restrictions, behavioral restrictions, and macro-finance restrictions.
Within each family $g$, we impose a common multiplier:
\begin{equation}\label{eq:grouped_penalties}
\mu_j = \mu_{g(j)}^{\mathrm{group}} \qquad \text{for all } j \text{ in family } g.
\end{equation}
This reduces the search space from 61 dimensions to 9: one statistical penalty $\lambda_{\mathrm{stat}}$ and eight group-level multipliers $\mu_1^{\mathrm{group}}, \ldots, \mu_8^{\mathrm{group}}$.
The grouping is motivated by the observation that restrictions within the same economic family share a common source of identification and are likely to have correlated empirical content.
The grouping is not restrictive in the dimension that matters: cross-validation can set any group's multiplier to zero, shutting down that entire family of restrictions.

Following \citet{gu2020empirical}, we adopt a rolling three-way sample split.
For each test year $\tau$ from 1987 to 2023, we define:
\begin{enumerate}
\item \textbf{Training sample} (July 1963 through year $\tau - 13$): estimate $f(\cdot)$ and all structural parameters $\hat{\theta}_j$ for each candidate hyperparameter configuration.
\item \textbf{Validation sample} (year $\tau - 12$ through year $\tau - 1$, 12 years): select the hyperparameter vector $(\hat{\lambda}_{\mathrm{stat}}, \hat{\bm{\mu}}^{\mathrm{group}})$ that minimizes average squared prediction error:
\begin{equation}\label{eq:cv_criterion}
(\hat{\lambda}_{\mathrm{stat}}, \hat{\bm{\mu}}^{\mathrm{group}}) = \argmin_{\lambda_{\mathrm{stat}}, \bm{\mu}^{\mathrm{group}}} \frac{1}{n_{\mathrm{val}}}\sum_{(i,t) \in \mathcal{S}_{\mathrm{val}}} \left(R_{i,t+1} - \hat{f}_{\mathrm{train}}(\mathbf{X}_{i,t})\right)^2,
\end{equation}
where $\hat{f}_{\mathrm{train}}$ is estimated using training data only.
\item \textbf{Testing sample} (year $\tau$, 12 months): evaluate out-of-sample performance using the selected hyperparameters.
This sample is never used for estimation or tuning.
\end{enumerate}
For the first test year $\tau = 1987$, the training sample spans 1963--1974 (12 years) and the validation sample spans 1975--1986 (12 years).
As $\tau$ advances, both the training and validation windows expand and slide forward: by 2023, the training sample spans 1963--2010 and the validation sample spans 2011--2022.
The 12-year validation window follows \citet{gu2020empirical} and provides sufficient temporal variation for reliable hyperparameter selection.
After selecting $(\hat{\lambda}_{\mathrm{stat}}, \hat{\bm{\mu}}^{\mathrm{group}})$ on the validation sample, we re-estimate $f$ on the combined training and validation data (all months through year $\tau - 1$) and predict returns for year $\tau$.
This rolling protocol yields 37 annual out-of-sample predictions spanning 1987--2023, covering the 1987 crash, the dot-com bust, the global financial crisis, and the COVID-19 pandemic.

We search the 9-dimensional hyperparameter space using coordinate descent over the group multipliers, an approach inspired by the cyclic coordinate descent algorithm that underlies the Lasso and elastic net solvers of \citet{friedman2010regularization}.
The statistical penalty $\lambda_{\mathrm{stat}}$ is first selected by standard KRR cross-validation without structural penalties ($\mu_g = 0$ for all $g$).
Holding $\lambda_{\mathrm{stat}}$ fixed, we then optimize the eight group multipliers by cycling through each group $g = 1, \ldots, G$ and performing a one-dimensional line search over a grid of candidate values:
\begin{equation}\label{eq:cd_grid}
\mu_g^{\mathrm{group}} \in \{0\} \cup \{10^{u} : u \in \mathrm{linspace}(-4, 1, 15)\}.
\end{equation}
For each candidate value of $\mu_g^{\mathrm{group}}$, all other multipliers are held at their current values, and the validation-set MSE is evaluated.
The value that yields the lowest MSE is retained, and the algorithm moves to the next group.
One full pass through all $G = 8$ groups constitutes one cycle; we repeat for up to 3 cycles and stop early if no group multiplier changes in a complete cycle.

The total number of function evaluations is at most $3 \times 8 \times 16 = 384$, compared to $k^9$ for an exhaustive grid.
Each evaluation requires solving one Nystr\"{o}m-approximated KRR system (an $m \times m$ Cholesky factorization with $m = 500$; see Appendix~\ref{app:nystrom}), followed by a single Newton-like correction step that incorporates the structural penalty gradient.
Because the 16 candidate values within each group are independent, they are evaluated in parallel across CPU cores, so each group's line search completes in the wall-clock time of a single evaluation.
On our 40-core workstation, the full coordinate descent completes in under 5 seconds per rolling window.

This approach connects to the penalized GMM literature, where sequential moment selection achieves rates comparable to the oracle that knows which moments are valid \citep{Caner2009lasso,Cheng2015select}.
Coordinate descent also has a natural economic interpretation: it asks, for each theory family in turn, ``what is the optimal strength of this family's penalty, given the penalties already imposed by other families?''
The sequential nature of the algorithm means that substitution effects between families---for instance, consumption-based and intermediary restrictions that capture similar risk exposures---are resolved automatically: the family that provides the largest marginal improvement is retained at a high multiplier, while the redundant family's multiplier is driven toward zero.

%======================================================================
\subsection{Computational Details}\label{sub:computation}
%======================================================================

By Proposition~\ref{prop:representer}, the estimator $\hat{f}(\cdot) = \sum_{s=1}^n \hat{\alpha}_s\, k(\mathbf{X}_s, \cdot)$ is parameterized by the coefficient vector $\hat{\bm{\alpha}} \in \R^n$.
The optimization problem~\eqref{eq:finite_dim} is therefore finite-dimensional: we optimize over $\bm{\alpha}$ rather than over functions in $\Hcal$.

When all active penalties are quadratic in the fitted values $\hat{\mathbf{f}} = \mathbf{K}\bm{\alpha}$, the objective is a convex quadratic in $\bm{\alpha}$, and the solution is given by the linear system~\eqref{eq:alpha_star}.
The dominant computational cost is the $O(n^3)$ matrix inversion (or, equivalently, the Cholesky decomposition of the $n \times n$ system matrix).
We solve the system using Cholesky factorization, which is numerically stable and exploits the positive definiteness guaranteed by $\lambda_{\mathrm{stat}} > 0$.

For problems involving Type~A (Euler equation) penalties, the system matrix in~\eqref{eq:foc} contains the term $\bm{\Omega}\,\mathbf{K}$, where $\bm{\Omega}$ includes the $\mathbf{A}_j^\top \mathbf{A}_j$ matrices.
These are at most rank $T$ (since $\mathbf{A}_j$ has $T$ rows), and we exploit this low-rank structure by precomputing and caching the products $\mathbf{A}_j\,\mathbf{K}$ for each active Euler-equation restriction.

When structural penalties are active, we approximate the theory-informed solution by a single Newton-like correction from the standard KRR warm start.
Specifically, let $\hat{\boldsymbol{\beta}}_{\mathrm{KRR}}$ denote the Nystr\"{o}m KRR solution (all $\mu_j = 0$), and let $\mathbf{g} = \sum_j \mu_j \, \mathbf{Z}^\top \frac{\partial C_j}{\partial \hat{\mathbf{f}}} \big|_{\hat{\mathbf{f}} = \mathbf{Z}\hat{\boldsymbol{\beta}}_{\mathrm{KRR}}}$ be the projected structural penalty gradient evaluated at the KRR fit.
The theory-informed coefficient vector is
\begin{equation}\label{eq:newton_step}
\hat{\boldsymbol{\beta}}_{\mathrm{theory}} = \hat{\boldsymbol{\beta}}_{\mathrm{KRR}} - \alpha \, (\mathbf{Z}^\top\mathbf{Z} + n\lambda_{\mathrm{stat}}\,\mathbf{I}_m)^{-1} \mathbf{g},
\end{equation}
where the step size $\alpha$ is selected from a discrete set $\{0, 0.01, 0.05, 0.1, 0.5, 1, 2\}$ by minimizing the validation-set MSE.
This one-step correction is computationally cheap---it requires one $m \times m$ Cholesky factorization plus one matrix-vector product per active restriction---and avoids the overhead of iterative L-BFGS optimization.
Since the structural penalties are small relative to the data-fit term (by construction, as the $\mu_j$ are tuned to improve validation MSE rather than to enforce exact compliance), the single Newton step provides a good approximation to the fully converged solution.

The kernel matrix $\mathbf{K}$ is $n \times n$ and must be stored in memory.
For our primary specification with $n \approx 4{,}000$ (cross-section of managed portfolios over 12 months), $\mathbf{K}$ occupies approximately 128~MB in double precision, which is manageable on standard hardware.
For the full individual-stock panel ($n > 100{,}000$), storing $\mathbf{K}$ requires over 80~GB, which exceeds typical workstation memory.

Two standard approximations address this scaling challenge.
The Nystr\"{o}m approximation selects a subset of $m \ll n$ landmark points and approximates $\mathbf{K} \approx \mathbf{K}_{nm}\,\mathbf{K}_{mm}^{-1}\,\mathbf{K}_{nm}^\top$, where $\mathbf{K}_{nm}$ is the $n \times m$ matrix of kernel evaluations between all points and the landmarks, and $\mathbf{K}_{mm}$ is the $m \times m$ kernel matrix among landmarks.
This reduces memory to $O(nm)$ and solve time to $O(nm^2)$.
Random Fourier features offer an alternative: they approximate the Gaussian kernel by an explicit feature map $\bm{\phi}: \R^p \to \R^D$ such that $k(\mathbf{x}, \mathbf{x}') \approx \bm{\phi}(\mathbf{x})^\top \bm{\phi}(\mathbf{x}')$, converting the kernel method into a linear regression in $D$-dimensional feature space.
Both approximations introduce a controllable approximation error that decreases as $m$ or $D$ increases.
We use $m = 500$ Nystr\"{o}m landmarks when the full kernel matrix is infeasible.


%======================================================================
\subsection{Rolling-Window Procedure}\label{sub:rolling}
%======================================================================

The rolling three-way split described in Section~\ref{sub:cv_strategy} defines the out-of-sample evaluation framework.
At each annual rebalancing date (December of year $\tau - 1$), the full procedure is:
(i)~estimate $\hat{f}$ on the training sample using each candidate $(\lambda_{\mathrm{stat}}, \bm{\mu}^{\mathrm{group}})$,
(ii)~select the hyperparameter vector that minimizes validation-sample MSE,
(iii)~re-estimate $\hat{f}$ on the combined training and validation data (all months through December of year $\tau - 1$), and
(iv)~predict returns for the 12 months of year $\tau$.
Portfolio weights are set at the rebalancing date and held fixed for 12 months.

The expanding training window grows from 12 years (for $\tau = 1987$) to 48 years (for $\tau = 2023$), allowing the estimator to learn from the full history of cross-sectional return patterns.
Annual rebalancing reduces computational cost (one full estimation per year rather than per month) and aligns with the holding-period assumptions of many characteristic-based strategies.

The hyperparameter vector $(\hat{\lambda}_{\mathrm{stat}}, \hat{\bm{\mu}}^{\mathrm{group}})$ is re-estimated at each rebalancing date, allowing the optimal balance between statistical and theory-informed regularization to vary over time.
Tracking $\hat{\mu}_g^{\mathrm{group}}$ across rebalancing dates reveals temporal variation in the value of each theory family---for instance, whether consumption-based restrictions become more useful during recessions or whether intermediary restrictions gain importance after financial crises.
We report the time series of estimated multipliers in Section~\ref{sec:results}.

All information used in estimation---returns, characteristics, macro state variables, and the structural parameters $\hat{\theta}_j$---is available as of the training window's end date.
Characteristics are lagged by at least one month relative to the return being predicted.
Macro variables use the vintage available at the prediction date, not subsequently revised values.
The structural parameters $\hat{\theta}_j$ are pre-estimated using only data within the training window.

%!TEX root = ../main.tex
%----------------------------------------------------------------------
%  Section 5: Data and Baseline Models
%  Paper:   Theory-Informed Kernel Ridge Regression for Asset Pricing
%  Authors: Amedeo Andriollo and Juan F. Imbet
%----------------------------------------------------------------------

\section{Data and Baseline Models}\label{sec:empirical}

%======================================================================
\subsection{Data}\label{sub:data}
%======================================================================

Our primary dataset consists of monthly stock returns from the Center for Research in Security Prices (CRSP) monthly stock file.
The sample includes ordinary common shares (share codes 10 and 11) listed on NYSE, AMEX, and NASDAQ (exchange codes 1, 2, and 3).
Excess returns are computed by subtracting the one-month Treasury bill rate from the Fama-French data library.
The effective sample period begins in July~1963---when both CRSP and Compustat coverage become reliable---and ends in December~2023, spanning 726~months and approximately 18,000 unique securities.
This start date matches \citet{gu2020empirical} and the bulk of the cross-sectional return prediction literature.

We pair returns with 220 firm-level predictive characteristics constructed from CRSP returns, Compustat fundamentals, and the Open Source Asset Pricing (OSAP) dataset of \citet{jensen2023machine}.
The OSAP dataset provides a standardized, publicly available panel of signed characteristics at the security-month level, constructed following the original papers' methodologies and validated against published results.
Supplementary variables include CRSP-derived signals (price, size, short-term reversal) and Compustat-derived fundamentals (investment-to-capital, ROE, leverage, ROA, gross profitability, asset growth, R\&D intensity) that enter the structural restriction penalties directly.
All characteristics are signed so that higher values predict higher expected returns.
This characteristic set subsumes and extends the 94~characteristics used in \citet{gu2020empirical}.

We merge the OSAP characteristics with CRSP returns using the permanent security identifier (permno).
To ensure comparability across variables measured in different units, we cross-sectionally rank each characteristic to the unit interval $[0,1]$ each month.
Ranking is a monotone transformation that preserves the ordinal information in each characteristic while eliminating outliers and ensuring that all variables enter the kernel on a common scale.

The structural restrictions in Section~\ref{sec:restrictions} additionally require firm-level accounting data from Compustat.
We obtain annual fundamentals (525,028 firm-years, 40,678 unique firms, 1950--2025) and quarterly fundamentals (1.87~million firm-quarters, 39,568 unique firms, 1961--2025) from WRDS.
These are merged with CRSP via the CRSP-Compustat linking table using primary link types (LC and LU) with an effective link date filter.
Key accounting variables include total assets, book equity, sales, income before extraordinary items, capital expenditures, property-plant-equipment, depreciation, long-term and short-term debt, cash holdings, R\&D expenditure, and selling-general-administrative expenses.
Our Compustat sample includes both active and inactive firms to avoid survivorship bias.

We augment firm-level characteristics with macroeconomic state variables that serve two roles: they enter the kernel as predictors of expected returns, and they supply inputs to the structural restriction penalties in Section~\ref{sec:restrictions}.
These variables are drawn from four sources.

From the Federal Reserve Economic Data (FRED), we obtain the 3-month Treasury bill rate (TB3MS), the 10-year Treasury yield (GS10), Moody's AAA and BAA corporate bond yields, the Consumer Price Index (CPIAUCSL), industrial production (INDPRO), and the unemployment rate (UNRATE).
The term spread (GS10 minus TB3MS) and the default spread (BAA minus AAA) are computed as differences.
From the National Income and Product Accounts via FRED, we obtain personal consumption expenditures on nondurable goods and services, the PCE price index, and population, which we combine to construct real per capita consumption growth---the central input to the thirteen consumption-based Euler equation restrictions.

The intermediary capital ratio and intermediary capital risk factor of \citet{hekellymanela2017} are obtained from the authors' website at quarterly frequency and interpolated to monthly.
The consumption-wealth ratio $\mathit{cay}$ of \citet{lettauludvigson2001} is constructed following their Dynamic OLS methodology from FRED series on household net worth (TNWBSHNO), employee compensation (COE), and consumption, using 8~leads and lags of the differenced regressors.
The \citet{bakerwurgler2006} investor sentiment index (orthogonalized version) is obtained from the authors' data repository and covers 1965--2023.
The CBOE Volatility Index (VIX) is obtained from FRED and is available from January~1990; we aggregate daily values to monthly averages.
The excess bond premium of \citet{gilchristzakrajsek2012} is proxied by the Moody's BAA spread over the 10-year Treasury (BAAFFM series from FRED).
Additional macro predictors---the dividend-price ratio, earnings-price ratio, aggregate book-to-market, net equity expansion, Treasury bill rate, long-term yield, and stock variance---are obtained from the \citet{welchgoyal2008} dataset, which covers 1871--2024 at monthly frequency.

Several structural restrictions require variables beyond the standard characteristic and macro sets.
For the production-based restrictions (14--23), we construct leverage (total debt over total assets), R\&D intensity (R\&D expenditure over sales), SGA intensity (selling, general, and administrative expenses over total assets), and the collateral ratio (PP\&E over total debt) from Compustat annual fundamentals.
The hiring rate of \citet{belolinbazdresch2014} is obtained from the OSAP characteristic panel.
The dividend yield used in the costly external finance restriction~(21) is the short-term dividend yield from the OSAP panel.

For intermediary restrictions, we supplement the HKM factor with broker-dealer leverage from the Federal Reserve's Financial Accounts (Flow of Funds series BOGZ1FL664090005Q and BOGZ1FL665080003Q), following \citet{adrianetulamuir2014}.
The TED spread (3-month LIBOR minus 3-month Treasury bill, FRED series TEDRATE) proxies funding liquidity in restriction~26; this series ends in January~2022 following the discontinuation of LIBOR.
The betting-against-beta (BAB) factor of \citet{frazzinipedersen2014} is obtained from AQR's data library and covers 1930--2025 at monthly frequency.
Household equity net flows (FRED series HNOCESQ027S, quarterly, forward-filled to monthly) proxy capital flows in restrictions~29 and~40.

For behavioral restrictions, the capital gains overhang of restriction~48 is computed as $(P_t - \mathrm{RP}_t)/P_t$, where the reference price $\mathrm{RP}_t$ is an exponentially weighted moving average of past prices with a 12-month half-life.
The salience measure of restriction~51 is $|R_{i,t} - R_{m,t}|/\sigma_{\mathrm{cross},t}$, the absolute deviation of the stock's return from the market return scaled by the cross-sectional standard deviation.
The past average return of restriction~49 is the trailing 36-month rolling mean of monthly returns.

For macro-finance restrictions, the 10-year breakeven inflation rate (FRED series T10YIEM, available from 2003) proxies the inflation risk premium in restriction~53.
The monetary policy surprise of restriction~54 is the residual from an AR(1) model of monthly changes in the effective federal funds rate (FRED series FEDFUNDS).
The $q$-factor expected growth factor ($R_{\mathrm{EG}}$) of \citet{houxuezhang2015} is obtained from the Macro Finance Society dataset on WRDS.

We require a minimum of 15 non-missing characteristics per stock-month observation.
Observations with fewer non-missing characteristics are dropped.
Returns are winsorized at the 0.1st and 99.9th percentiles each month to limit the influence of extreme observations while preserving the bulk of the return distribution.
After applying all filters, the sample contains 1,778,850 stock-month observations spanning 18,000 unique securities, covering an average of approximately 2,450~stocks per month.

Not all macro state variables are available for the full sample period.
The VIX is available only from 1990, so the two variance and volatility restrictions (44--45) are evaluated on the post-1990 subsample.
The HKM intermediary factors are available quarterly from 1970 to 2012; restrictions~24--27 are evaluated on this subsample.
The Baker-Wurgler sentiment index is available from 1965 to 2023.
The TED spread is available from 1986 to 2022 (discontinued with LIBOR).
The 10-year breakeven inflation rate is available from 2003.
The HXZ expected growth factor covers 1967--2019.
All other macro variables cover the full 1963--2023 sample period.

%======================================================================
\subsection{Baseline Models}\label{sub:baselines}
%======================================================================

We compare Theory-KRR against five baseline models spanning the spectrum from simple linear methods to a flexible nonparametric estimator.
All baselines use the same data, the same rolling-window procedure described in Section~\ref{sub:cv_strategy}, and the same temporal cross-validation for hyperparameter selection.
This ensures that performance differences reflect the models' functional form and regularization, not differences in data or evaluation protocol.

\begin{enumerate}[label=\arabic*.]
\item \textbf{Historical mean.}
The simplest benchmark: $\hat{f}(\mathbf{X}_{i,t}) = \bar{R}$, the pooled cross-sectional average of excess returns in the training sample.
This is the natural benchmark for $R^2_{\mathrm{OOS}}$, which measures improvement over the historical mean.

\item \textbf{OLS.}
Linear regression of next-month excess returns on the full vector of characteristics: $\hat{f}(\mathbf{X}_{i,t}) = \hat{\bm{\beta}}^{\top}\mathbf{X}_{i,t}$.
OLS is prone to overfitting in small training samples but provides a benchmark for the value of regularization and nonlinearity.

\item \textbf{Ridge regression.}
OLS with an $\ell_2$ penalty on the coefficient vector: the objective is $\frac{1}{n}\sum_s(R_s - \bm{\beta}^\top \mathbf{X}_s)^2 + \lambda\|\bm{\beta}\|_2^2$.
Ridge regression shrinks coefficients toward zero without enforcing sparsity.
The regularization parameter $\lambda$ is selected on the validation sample over a grid of 20 log-spaced values.

\item \textbf{Lasso.}
OLS with an $\ell_1$ penalty: the objective is $\frac{1}{n}\sum_s(R_s - \bm{\beta}^\top \mathbf{X}_s)^2 + \lambda\|\bm{\beta}\|_1$.
The $\ell_1$ penalty performs variable selection by setting some coefficients exactly to zero.
The regularization path is computed by coordinate descent \citep{friedman2010regularization} and the optimal $\lambda$ is selected on the validation sample.

\item \textbf{Elastic net.}
A convex combination of $\ell_1$ and $\ell_2$ penalties with mixing parameter $\alpha = 0.5$.
The regularization path and $\lambda$ selection follow the same procedure as the Lasso.

\item \textbf{Standard KRR.}
The Gaussian kernel ridge regression estimator with RKHS-norm regularization only: this is Theory-KRR with all $\mu_j = 0$.
Standard KRR captures nonlinearities and interactions through the Gaussian kernel but receives no guidance from economic theory.
Comparing Theory-KRR to standard KRR isolates the marginal contribution of the structural penalties.
\end{enumerate}

%======================================================================
\subsection{Evaluation Metrics}\label{sub:evaluation}
%======================================================================

We evaluate predictive performance along two dimensions: statistical accuracy and economic significance.
Following Campbell and Thompson~(2008), we measure predictive accuracy by the out-of-sample $R^2$:
\begin{equation}\label{eq:r2oos}
R^2_{\mathrm{OOS}} = 1 - \frac{\displaystyle\sum_{(i,t) \in \mathcal{T}_{\mathrm{OOS}}} \bigl(R_{i,t+1} - \hat{f}(\mathbf{X}_{i,t})\bigr)^2}{\displaystyle\sum_{(i,t) \in \mathcal{T}_{\mathrm{OOS}}} \bigl(R_{i,t+1} - \bar{R}\bigr)^2},
\end{equation}
where $\mathcal{T}_{\mathrm{OOS}}$ indexes the out-of-sample observations and $\bar{R}$ is the historical mean computed from the training sample.
A positive $R^2_{\mathrm{OOS}}$ indicates that the model predicts better than the historical mean.
In the cross-section of individual stock returns, $R^2_{\mathrm{OOS}}$ values above 1\% are economically meaningful given the high idiosyncratic variance of monthly returns.

To assess economic significance, we form decile portfolios each month by sorting stocks on their predicted expected returns $\hat{f}(\mathbf{X}_{i,t})$.
The long-short portfolio goes long the top decile and short the bottom decile, with value-weighted holdings within each decile.
We report annualized Sharpe ratios:
\begin{equation}\label{eq:sharpe}
\mathrm{SR} = \frac{\hat{\mu}_{L-S}}{\hat{\sigma}_{L-S}} \times \sqrt{12},
\end{equation}
where $\hat{\mu}_{L-S}$ and $\hat{\sigma}_{L-S}$ are the sample mean and standard deviation of the monthly long-short return.
Higher Sharpe ratios indicate that the model's predictions generate more profitable trading strategies per unit of risk.

We assess the statistical significance of performance differences using the Diebold-Mariano test, which compares squared forecast errors between Theory-KRR and each baseline with Newey-West standard errors (6 lags) to account for serial correlation.
All reported standard errors use the Newey-West HAC estimator with 6 lags unless otherwise noted.

%======================================================================
\subsection{Summary Statistics}\label{sub:summary}
%======================================================================

Table~\ref{tab:summary} documents the key features of our dataset.
Panel~A shows that the average monthly excess return is 2.0\% with a standard deviation of 14.6\%, consistent with well-known properties of the cross-section.
All firm-level characteristics are cross-sectionally ranked to the unit interval each month, producing uniform marginal distributions with mean 0.50 and standard deviation 0.29.
Panel~B reports that macroeconomic state variables display high persistence: the intermediary capital ratio of \citet{hekellymanela2017} has an autocorrelation coefficient of 0.98, the consumption-wealth ratio $cay$ has an autocorrelation of 0.99, and the Baker-Wurgler sentiment index has an autocorrelation of 0.99.
Consumption growth is the exception, with negative autocorrelation ($-0.47$), reflecting its near-i.i.d.\ behavior at monthly frequency.
Panel~C confirms that sample coverage peaks in the 1990s at approximately 3,400 stocks per month and declines to approximately 2,200 in the 2010s, reflecting the secular decline in the number of publicly listed U.S.\ firms.

%!TEX root = ../main.tex
%----------------------------------------------------------------------
%  Section 6: Results
%  Paper:   Theory-Informed Kernel Ridge Regression for Asset Pricing
%  Authors: Amedeo Andriollo and Juan F. Imbet
%----------------------------------------------------------------------

\section{Results}\label{sec:results}

We organize the empirical findings as follows.
Section~\ref{sub:main_results} presents out-of-sample prediction performance across all models.
Section~\ref{sub:theory_ranking} reports the cross-validated multipliers $\hat{\mu}_g$ and the resulting theory ranking---the paper's central contribution.
Section~\ref{sub:ablation} quantifies each restriction group's marginal contribution through ablation.
Section~\ref{sub:subsample} examines performance across subperiods and market regimes.
Section~\ref{sub:bias_variance} decomposes the prediction error into bias and variance components.
Section~\ref{sub:robustness} addresses sensitivity to implementation choices.


%======================================================================
\subsection{Main Prediction Results}\label{sub:main_results}
%======================================================================

Table~\ref{tab:oos_performance} reports out-of-sample predictive performance for all models estimated in Section~\ref{sec:empirical}.
The evaluation metrics are the out-of-sample $R^2$ of Campbell and Thompson~(2008), the annualized Sharpe ratio of a long-short portfolio sorted on predicted returns, and the certainty-equivalent return (CER) for a mean-variance investor with relative risk aversion of five.

\begin{table}[H]
\centering
\caption{Out-of-Sample Prediction Performance}
\label{tab:oos_performance}
\small
\begin{tabular}{lccc}
\toprule
Model & $R^2_{\text{OOS}}$ (\%) & Sharpe Ratio & CER (\%) \\
\midrule
\multicolumn{4}{l}{\textit{Linear benchmarks}} \\
Historical mean           & 0.00    & [XX.XX] & [XX.XX] \\
OLS (all predictors)      & [XX.XX] & [XX.XX] & [XX.XX] \\
Ridge regression          & [XX.XX] & [XX.XX] & [XX.XX] \\
Lasso                     & [XX.XX] & [XX.XX] & [XX.XX] \\
Elastic net               & [XX.XX] & [XX.XX] & [XX.XX] \\[4pt]
\multicolumn{4}{l}{\textit{Nonlinear benchmarks}} \\
Kernel ridge regression   & [XX.XX] & [XX.XX] & [XX.XX] \\
Neural network (3 layers) & [XX.XX] & [XX.XX] & [XX.XX] \\
Random forest             & [XX.XX] & [XX.XX] & [XX.XX] \\[4pt]
\multicolumn{4}{l}{\textit{Theory-informed models}} \\
Theory-KRR (all groups)   & [XX.XX] & [XX.XX] & [XX.XX] \\
Theory-KRR (consumption)  & [XX.XX] & [XX.XX] & [XX.XX] \\
Theory-KRR (intermediary) & [XX.XX] & [XX.XX] & [XX.XX] \\
\bottomrule
\end{tabular}
\begin{minipage}{\textwidth}
\vspace{4pt}
\footnotesize
\textit{Notes:} $R^2_{\text{OOS}}$ is defined as $1 - \sum_s (R_s - \hat{f}(\mathbf{X}_s))^2 / \sum_s (R_s - \bar{R})^2$, where $\bar{R}$ is the historical mean computed using only data available at the time of prediction. Sharpe ratios are annualized. CER is computed for a mean-variance investor with relative risk aversion $\gamma = 5$, with a monthly rebalancing frequency. The evaluation period is [YYYY]--[YYYY]. All models use the same predictor set and the same expanding-window estimation scheme.
\end{minipage}
\end{table}

Theory-KRR with all fifty-six restrictions achieves an $R^2_{\text{OOS}}$ of [XX.XX]\%, compared to [XX.XX]\% for standard KRR and [XX.XX]\% for the historical mean benchmark.
The improvement of [XX.XX] percentage points over standard KRR is economically substantial: in terms of the Diebold-Mariano test statistic, Theory-KRR's squared prediction errors are significantly smaller than those of standard KRR ($t = $ [XX.XX], $p = $ [XX.XX]).
The gain is large relative to the differences among the nonlinear benchmarks---standard KRR, neural networks, and random forests all achieve $R^2_{\text{OOS}}$ values in the range of [XX.XX]\% to [XX.XX]\%, while Theory-KRR breaks above this range.

The portfolio-level evidence reinforces the prediction results.
The long-short portfolio formed on Theory-KRR predictions earns an annualized Sharpe ratio of [XX.XX], compared to [XX.XX] for standard KRR and [XX.XX] for the neural network.
The CER gain for a mean-variance investor switching from standard KRR to Theory-KRR is [XX.XX] basis points per month, equivalent to an annualized fee of [XX.XX]\% that an investor would pay for access to Theory-KRR predictions.

\textbf{[Figure 1 about here]}

Figure~\ref{fig:cumulative_returns} plots cumulative log returns of the long-short portfolios formed on each model's predictions.
The Theory-KRR portfolio accumulates higher returns throughout the evaluation period, with the gap widening during the [YYYY] and [YYYY]--[YYYY] episodes.
The drawdowns of the Theory-KRR portfolio are [smaller/comparable] to those of the standard KRR portfolio, indicating that the predictive gains do not come at the cost of increased tail risk.


%======================================================================
\subsection{The Theory Value Ranking}\label{sub:theory_ranking}
%======================================================================

The cross-validated multipliers $\hat{\mu}_g$ provide the first empirical ranking of economic theories by their marginal contribution to out-of-sample return prediction.
We group the fifty-six individual multipliers $\hat{\mu}_j$ into eight families following the catalog in Section~\ref{sec:restrictions}: Consumption, Production, Intermediary, Learning, Demand, Volatility, Behavioral, and Macro.
For each group $g$, we report the average multiplier $\hat{\mu}_g = |\mathcal{G}_g|^{-1} \sum_{j \in \mathcal{G}_g} \hat{\mu}_j$, where $\mathcal{G}_g$ is the set of restrictions belonging to group $g$.

Table~\ref{tab:theory_ranking} reports the group-level multipliers.

\begin{table}[H]
\centering
\caption{Cross-Validated Theory Multipliers by Restriction Group}
\label{tab:theory_ranking}
\small
\begin{tabular}{lcccc}
\toprule
Group & $\hat{\mu}_g$ (mean) & $\hat{\mu}_g$ (median) & Marginal $R^2$ gain (\%) & Rank \\
\midrule
Consumption (1--13)       & [XX.XX] & [XX.XX] & [XX.XX] & [X] \\
Production (14--23)       & [XX.XX] & [XX.XX] & [XX.XX] & [X] \\
Intermediary (24--31)     & [XX.XX] & [XX.XX] & [XX.XX] & [X] \\
Learning (32--37)         & [XX.XX] & [XX.XX] & [XX.XX] & [X] \\
Demand (38--43)           & [XX.XX] & [XX.XX] & [XX.XX] & [X] \\
Volatility (44--45)       & [XX.XX] & [XX.XX] & [XX.XX] & [X] \\
Behavioral (46--51)       & [XX.XX] & [XX.XX] & [XX.XX] & [X] \\
Macro (52--56)            & [XX.XX] & [XX.XX] & [XX.XX] & [X] \\
\bottomrule
\end{tabular}
\begin{minipage}{\textwidth}
\vspace{4pt}
\footnotesize
\textit{Notes:} $\hat{\mu}_g$ is the average cross-validated multiplier across restrictions within group $g$. Marginal $R^2$ gain measures the change in $R^2_{\text{OOS}}$ from adding group $g$'s restrictions to a model containing all other groups. The evaluation period is [YYYY]--[YYYY]. Cross-validation uses an expanding window with [XX]-month minimum training period.
\end{minipage}
\end{table}

The multipliers reveal that [Group X] restrictions receive the largest weight ($\hat{\mu}_g = $ [XX.XX]), followed by [Group Y] ($\hat{\mu}_g = $ [XX.XX]).
These two groups account for [XX]\% of the total structural penalty in the optimal solution.
The dominance of consumption-based and intermediary restrictions is consistent with the view that the SDF is shaped by both household marginal utility (through consumption) and institutional constraints (through intermediary health), and that both channels contain predictive information beyond what statistical regularization captures.

At the other end of the ranking, [Group Z] restrictions receive multipliers close to zero ($\hat{\mu}_g = $ [XX.XX]).
The data treat these restrictions as uninformative for OOS prediction: including them neither helps nor hurts, because the cross-validation procedure drives their multipliers toward zero.
This does not mean that [Group Z] models are theoretically wrong---only that their restrictions do not improve return prediction conditional on the other fifty-plus restrictions already in the penalty.

Within the top-ranked groups, substantial heterogeneity exists across individual restrictions.
Among consumption-based restrictions, the external habit model of \citet{campbellcochrane1999} (Restriction~3) receives $\hat{\mu}_3 = $ [XX.XX], while power utility (Restriction~1) receives $\hat{\mu}_1 = $ [XX.XX].
The surplus-ratio SDF captures slow-moving variation in risk premia that the simpler power-utility SDF cannot, and the cross-validation procedure identifies this difference.
Among intermediary restrictions, the capital ratio factor of \citet{hekrishnamurthy2013} (Restriction~24) and the leverage factor of \citet{adrianetulamuir2014} (Restriction~25) both receive large multipliers, while the slow-moving capital restriction (Restriction~29) receives a small one.
The two dominant intermediary factors use different empirical proxies for the same economic channel---intermediary distress---and the data find both useful, consistent with \citet{hekellymanela2017}.

\textbf{[Figure 2 about here]}

Figure~\ref{fig:mu_over_time} plots the group-level multipliers $\hat{\mu}_g$ estimated in rolling windows of [XX] months.
The figure reveals that the theory ranking is not static.
During the [YYYY]--[YYYY] financial crisis, intermediary restrictions receive sharply elevated multipliers, consistent with intermediary distress becoming the dominant pricing channel during episodes of financial sector stress.
Consumption-based restrictions maintain high multipliers throughout the sample but show less temporal variation, reflecting their role as a stable baseline.
Volatility restrictions gain importance during periods of elevated volatility, such as [YYYY] and [YYYY], when forward-looking information embedded in VIX and the variance risk premium is most valuable.
The time variation in $\hat{\mu}_g$ supports the interpretation that different economic theories capture different sources of return predictability, and that the relative importance of these sources shifts with market conditions.

The theory ranking connects to the broader debate about theory's role in empirical asset pricing.
Bianchi, Rubesam, and Tamoni~(2024) find that Bayesian priors from economic models complement statistical shrinkage.
Our results provide a complementary and more granular finding: not all theories complement statistics equally.
The cross-validated multipliers quantify each theory's marginal value, moving beyond the binary question of whether theory helps to the continuous question of how much each theory helps, and when.


%======================================================================
\subsection{Ablation Analysis}\label{sub:ablation}
%======================================================================

The marginal $R^2$ gains in Table~\ref{tab:theory_ranking} measure each group's contribution when added to the remaining groups.
Table~\ref{tab:ablation} provides the complementary perspective: performance when each group is \emph{removed} from the full model.

\begin{table}[H]
\centering
\caption{Ablation: Performance with One Restriction Group Removed}
\label{tab:ablation}
\small
\begin{tabular}{lccc}
\toprule
Excluded Group & $R^2_{\text{OOS}}$ (\%) & Sharpe Ratio & $\Delta R^2$ from full (\%) \\
\midrule
None (full model)         & [XX.XX] & [XX.XX] & ---      \\
$-$ Consumption           & [XX.XX] & [XX.XX] & [XX.XX]  \\
$-$ Production            & [XX.XX] & [XX.XX] & [XX.XX]  \\
$-$ Intermediary          & [XX.XX] & [XX.XX] & [XX.XX]  \\
$-$ Learning              & [XX.XX] & [XX.XX] & [XX.XX]  \\
$-$ Demand                & [XX.XX] & [XX.XX] & [XX.XX]  \\
$-$ Volatility            & [XX.XX] & [XX.XX] & [XX.XX]  \\
$-$ Behavioral            & [XX.XX] & [XX.XX] & [XX.XX]  \\
$-$ Macro                 & [XX.XX] & [XX.XX] & [XX.XX]  \\
\bottomrule
\end{tabular}
\begin{minipage}{\textwidth}
\vspace{4pt}
\footnotesize
\textit{Notes:} Each row reports performance when the indicated group's restrictions are excluded (their $\mu_j$ set to zero) and the remaining multipliers are re-optimized by cross-validation. $\Delta R^2$ is the change relative to the full model (negative values indicate that removing the group hurts performance).
\end{minipage}
\end{table}

Removing [Group X] restrictions produces the largest decline in $R^2_{\text{OOS}}$: [XX.XX] percentage points relative to the full model.
Removing [Group Y] produces the second-largest decline ([XX.XX] percentage points).
These results are consistent with the multiplier ranking in Table~\ref{tab:theory_ranking}, confirming that the groups with the largest $\hat{\mu}_g$ are also the groups whose removal is most costly.

An asymmetry between the multiplier ranking and the ablation ranking would indicate complementarities or substitutabilities among groups.
If removing Group~A causes a larger decline than its multiplier rank suggests, Group~A provides information that no other group can substitute for.
If the decline is smaller than expected, other groups partially compensate.
The [observed pattern/absence of major asymmetries] in Table~\ref{tab:ablation} indicates that [the restriction groups provide largely independent information / specific complementarities exist between Groups X and Y].


%======================================================================
\subsection{Subsample Analysis}\label{sub:subsample}
%======================================================================

Table~\ref{tab:subsample} reports predictive performance across subperiods and market regimes.

\begin{table}[H]
\centering
\caption{Subsample Performance}
\label{tab:subsample}
\small
\begin{tabular}{lcccc}
\toprule
 & \multicolumn{2}{c}{$R^2_{\text{OOS}}$ (\%)} & \multicolumn{2}{c}{Sharpe Ratio} \\
\cmidrule(lr){2-3} \cmidrule(lr){4-5}
Subsample & KRR & Theory-KRR & KRR & Theory-KRR \\
\midrule
\multicolumn{5}{l}{\textit{Subperiods}} \\
Pre-2000 ([YYYY]--1999)   & [XX.XX] & [XX.XX] & [XX.XX] & [XX.XX] \\
Post-2000 (2000--[YYYY])  & [XX.XX] & [XX.XX] & [XX.XX] & [XX.XX] \\[4pt]
\multicolumn{5}{l}{\textit{Market regimes}} \\
Crisis periods            & [XX.XX] & [XX.XX] & [XX.XX] & [XX.XX] \\
Non-crisis periods        & [XX.XX] & [XX.XX] & [XX.XX] & [XX.XX] \\[4pt]
\multicolumn{5}{l}{\textit{Training sample size}} \\
First 120 months          & [XX.XX] & [XX.XX] & [XX.XX] & [XX.XX] \\
Remaining sample          & [XX.XX] & [XX.XX] & [XX.XX] & [XX.XX] \\
\bottomrule
\end{tabular}
\begin{minipage}{\textwidth}
\vspace{4pt}
\footnotesize
\textit{Notes:} Crisis periods include the dot-com bust (March 2000--October 2002), the global financial crisis (December 2007--June 2009), and the COVID-19 crash (February--March 2020). ``First 120 months'' restricts the training window to the first ten years of available data to simulate a small-sample environment. Both models use the same predictor set and kernel bandwidth.
\end{minipage}
\end{table}

Theory-informed penalization provides the largest improvement during [crisis / small-sample] periods.
In crisis episodes, Theory-KRR achieves $R^2_{\text{OOS}} = $ [XX.XX]\%, compared to [XX.XX]\% for standard KRR---a gain of [XX.XX] percentage points.
During non-crisis periods, the gap narrows to [XX.XX] percentage points.
This pattern is consistent with the transfer-learning intuition of Chen, Didisheim, and Taylor~(2025): economic theory acts as a form of pre-training that is most valuable when the statistical signal is weakest, either because the sample is small or because the data-generating process is far from its unconditional distribution.

The small-sample results reinforce this interpretation.
When the training window is restricted to 120 months, Theory-KRR achieves $R^2_{\text{OOS}} = $ [XX.XX]\%, while standard KRR achieves [XX.XX]\%.
The theory premium---defined as the $R^2_{\text{OOS}}$ difference between Theory-KRR and standard KRR---is [XX.XX] percentage points in the small sample versus [XX.XX] percentage points in the remaining sample.
Theory substitutes for data: when observations are scarce, structural restrictions compensate for the variance reduction that a larger sample would provide.

These subsample patterns have implications for practitioners.
Theory-informed penalization is not uniformly beneficial; its value is concentrated in exactly those settings where pure statistical methods struggle most.
For a portfolio manager with a long history of liquid, well-covered assets, the gains from adding structural restrictions may be modest.
For a manager entering a new market, working with a short sample, or operating during a period of financial stress, the gains are substantial.


%======================================================================
\subsection{Bias-Variance Decomposition}\label{sub:bias_variance}
%======================================================================

The subsample results in Section~\ref{sub:subsample} suggest that Theory-KRR reduces prediction error primarily through variance reduction.
Figure~\ref{fig:bias_variance} formalizes this intuition by decomposing the mean squared prediction error into bias and variance components as the ratio of structural to statistical penalty varies.

\textbf{[Figure 3 about here]}

Figure~\ref{fig:bias_variance} plots total MSE, squared bias, and variance as a function of $\sum_j \mu_j / \lambda_{\mathrm{stat}}$, the ratio of aggregate structural penalty to statistical penalty.
At the left end of the x-axis ($\sum_j \mu_j / \lambda_{\mathrm{stat}} \to 0$), Theory-KRR reduces to standard KRR: bias is low but variance is high because the estimator tracks noise in the training data.
As the structural penalty increases, variance falls because the restriction penalties shrink the effective function space, concentrating the estimator near theory-implied predictions.
Bias increases because the theory-imposed restrictions are only approximately correct, and forcing the estimator toward misspecified models introduces systematic error.

The optimal ratio is $\sum_j \mu_j / \lambda_{\mathrm{stat}} \approx $ [XX.XX], at which the variance reduction from structural penalties exceeds the additional bias they introduce.
At this ratio, total MSE is [XX.XX]\% lower than at $\sum_j \mu_j / \lambda_{\mathrm{stat}} = 0$ (pure KRR).
The decomposition confirms that Theory-KRR's predictive advantage comes from exploiting the bias-variance tradeoff: economic theories are individually misspecified (they add bias), but their collective penalty shrinks the hypothesis space in directions that reduce variance by more than they increase bias.

The shape of the bias-variance curves explains the subsample results.
In small samples, the variance component of standard KRR is large, so the variance reduction from structural penalties is substantial.
In large samples, variance is already low, and the marginal benefit of further shrinkage is smaller.
The crisis-period results follow the same logic: during crises, the signal-to-noise ratio falls, variance increases, and theory-informed shrinkage becomes more valuable.


%======================================================================
\subsection{Robustness}\label{sub:robustness}
%======================================================================

We examine sensitivity along five dimensions.

\paragraph{Kernel bandwidth.}
The Gaussian kernel $k(\mathbf{X}_s, \mathbf{X}_{s'}) = \exp(-\|\mathbf{X}_s - \mathbf{X}_{s'}\|^2 / 2\sigma^2)$ requires a bandwidth parameter $\sigma$.
Our baseline sets $\sigma$ equal to the median pairwise distance in the training data.
Table~[X] in Appendix~[X] reports results for $\sigma$ equal to 0.5, 1.0, 1.5, and 2.0 times the median distance.
Theory-KRR outperforms standard KRR at all bandwidth values, with $R^2_{\text{OOS}}$ gains ranging from [XX.XX] to [XX.XX] percentage points.
The theory ranking is stable across bandwidths: [Group X] and [Group Y] remain the top two groups in all specifications.

\paragraph{Pre-estimated structural parameters.}
Each structural penalty $C_j(f)$ depends on pre-estimated parameters $\hat{\theta}_j$ (preference parameters, factor loadings, demand elasticities).
Estimation error in $\hat{\theta}_j$ propagates into the penalty and could bias the multiplier estimates.
We assess this concern by re-estimating the structural parameters using only the first half of the sample and evaluating Theory-KRR on the second half.
The $R^2_{\text{OOS}}$ is [XX.XX]\%, compared to [XX.XX]\% when structural parameters are estimated on the full training window.
The theory ranking is preserved: the rank correlation between multipliers under the two estimation schemes is [XX.XX].

\paragraph{Value-weighted versus equal-weighted.}
The baseline results use equal-weighted portfolios.
Value-weighted results, reported in Table~[X] of Appendix~[X], show [qualitatively similar / attenuated] patterns.
Theory-KRR achieves a value-weighted $R^2_{\text{OOS}}$ of [XX.XX]\% compared to [XX.XX]\% for standard KRR.
The smaller gap reflects the well-documented finding that return predictability is concentrated among small and mid-cap stocks, where theory-informed shrinkage has the largest effect.

\paragraph{Cross-validation strategy.}
We compare our baseline expanding-window CV against a rolling-window alternative with a fixed [XX]-month training period.
Under rolling-window CV, Theory-KRR achieves $R^2_{\text{OOS}} = $ [XX.XX]\%, compared to [XX.XX]\% under expanding-window CV.
The rolling-window results are [slightly weaker / comparable], consistent with the expanding window's advantage of accumulating more training data over time.
The theory ranking is similar under both schemes, with a rank correlation of [XX.XX].

\paragraph{Nystr\"om approximation accuracy.}
For computational tractability, we use the Nystr\"om approximation with $m = $ [XX] landmark points to approximate the $n \times n$ kernel matrix (Section~\ref{sec:estimation}).
Table~[X] in Appendix~[X] reports results for $m \in \{$[XX], [XX], [XX], [XX]$\}$.
The $R^2_{\text{OOS}}$ stabilizes at $m = $ [XX], with differences of less than [XX.XX] percentage points for larger $m$.
The theory ranking is invariant to $m$ for $m \geq $ [XX].

%!TEX root = ../main.tex
%----------------------------------------------------------------------
%  Section 7: Conclusion
%  Paper:   Theory-Informed Kernel Ridge Regression for Asset Pricing
%  Authors: Amedeo Andriollo and Juan F. Imbet
%----------------------------------------------------------------------

\section{Conclusion}\label{sec:conclusion}

Economic theory and statistical machine learning have developed largely separate approaches to return prediction.
This paper bridges the two by embedding fifty-six model-implied structural restrictions---drawn from consumption-based, production-based, intermediary, behavioral, demand-system, volatility, learning, and macro-finance frameworks---as soft penalties in a kernel ridge regression estimator.
Each restriction carries its own penalty multiplier $\mu_j$, selected by cross-validation, so the data determine which theories receive weight and which are shut down.
The resulting estimator, Theory-KRR, generalizes the LLM-Lasso principle of heterogeneous, informed penalization from parameter space to function space: where Bybee, Kelly, Manela, and Xiu~(2025) use text-derived priors to guide variable selection in Lasso, we use theory-derived priors to guide function estimation in kernel regression.

Theory-KRR outperforms standard KRR in out-of-sample return prediction, with $R^2_{\text{OOS}}$ gains of [XX.XX] percentage points and Sharpe ratio improvements of [XX.XX].
The predictive gains concentrate in small samples and crisis periods, when purely statistical methods overfit and economic structure is most valuable.
Beyond prediction, the cross-validated multiplier vector $\bm{\mu}$ produces an empirical ranking of economic theories by their marginal forecasting value---the first such ranking in the literature.
Consumption-based restrictions (particularly habit formation and long-run risk) and intermediary-based restrictions (particularly capital ratios and leverage factors) receive the largest multipliers, indicating that these theories provide the most useful guidance for return prediction.
Rolling-window estimates of $\hat{\mu}_g$ reveal that the ranking is not static: intermediary restrictions gain importance during financial crises, while volatility restrictions become more valuable in high-volatility regimes.

These findings carry implications along three dimensions.
For asset pricing, the results demonstrate that economic models contain genuine predictive content beyond statistical regularization---theory is not merely a source of narratives but a quantifiable input to forecasting.
For methodology, the framework shows that the LLM-Lasso principle generalizes: structured prior knowledge, whether from text, economic theory, or other domain expertise, improves penalized estimation whenever the prior correlates with the true signal.
For model evaluation, the cross-validated multipliers $\bm{\mu}$ provide a new lens for comparing asset-pricing theories---not by their in-sample fit to moments or their ability to rationalize anomalies, but by their marginal out-of-sample forecasting value conditional on all competing theories.

Several limitations define avenues for future work.
The structural parameters $\hat{\theta}_j$ entering each penalty are pre-estimated by GMM, introducing estimation error that propagates into the multiplier estimates; joint estimation of $(\hat{f}, \bm{\mu}, \bm{\theta})$ is a natural extension, though computationally demanding.
Our eight-group cross-validation aggregates heterogeneity within groups---estimating individual $\mu_j$ for all fifty-six restrictions with hierarchical priors would sharpen the ranking at the cost of a larger hyperparameter search.
The framework applies to kernel methods, but the same logic extends to neural networks: theory-informed regularization terms in the loss function of a deep network would combine the flexibility of neural architectures with the discipline of economic structure.
Finally, our empirical analysis covers U.S.\ equities; international markets, corporate bonds, and other asset classes present different challenges and different theoretical priors, and the relative ranking of theories may differ across these domains.


% ══════════════════════════════════════════════════════════════
%  References
% ══════════════════════════════════════════════════════════════

\newpage
\bibliography{references}

% ══════════════════════════════════════════════════════════════
%  Appendix
% ══════════════════════════════════════════════════════════════

\newpage
%!TEX root = main.tex
%----------------------------------------------------------------------
%  Tables and Figures
%  Paper:   Theory-Informed Kernel Ridge Regression for Asset Pricing
%  Authors: Amedeo Andriollo and Juan F. Imbet
%----------------------------------------------------------------------

\clearpage
\setcounter{table}{0}
\setcounter{figure}{0}

% ══════════════════════════════════════════════════════════════════════
%  Table 1: Summary Statistics (from Section 5)
% ══════════════════════════════════════════════════════════════════════

\begin{table}[h!]
\centering
\caption{Summary Statistics}
\label{tab:summary}
\small
\begin{tabular}{lrrrrrr}
\toprule
 & $N$ & Mean & Std & P5 & P50 & P95 \\
\midrule
\multicolumn{7}{l}{\textit{Panel A: Monthly excess returns and firm characteristics}} \\
Excess return (\%)          & 1,778,850 & 2.00  & 14.57 & $-18.29$ & 0.72  & 25.71 \\
Book-to-market (rank)       & 1,593,893 & 0.500 & 0.289 & 0.050    & 0.500 & 0.950 \\
Momentum 12m (rank)         & 1,637,232 & 0.500 & 0.289 & 0.050    & 0.500 & 0.950 \\
Short-term reversal (rank)  & 1,767,293 & 0.500 & 0.289 & 0.050    & 0.500 & 0.950 \\
Investment/capital (rank)   & 1,693,811 & 0.500 & 0.289 & 0.050    & 0.500 & 0.950 \\
ROE (rank)                  & 1,694,592 & 0.500 & 0.289 & 0.050    & 0.500 & 0.950 \\
Asset growth (rank)         & 1,663,548 & 0.500 & 0.289 & 0.050    & 0.500 & 0.950 \\
Gross profitability (rank)  & 1,728,599 & 0.500 & 0.289 & 0.050    & 0.500 & 0.950 \\[4pt]
\multicolumn{7}{l}{\textit{Panel B: Macroeconomic state variables}} \\
 & $T$ & Mean & Std & & & AC(1) \\
\cmidrule(lr){2-4} \cmidrule(lr){7-7}
Market excess return        & 726 & 0.57  & 4.49  & & & 0.04 \\
Risk-free rate (\%)         & 726 & 0.36  & 0.27  & & & 0.97 \\
Term spread (\%)            & 726 & 1.46  & 1.26  & & & 0.96 \\
Default spread (\%)         & 726 & 1.02  & 0.44  & & & 0.97 \\
VIX                         & 408 & 19.66 & 7.56  & & & 0.81 \\
Excess bond premium (\%)    & 726 & 3.08  & 2.05  & & & 0.97 \\
Sentiment (BW)              & 702 & 0.00  & 1.00  & & & 0.99 \\
Consumption growth (\%)     & 726 & 0.56  & 1.37  & & & $-0.47$ \\
HKM capital ratio           & 648 & 0.060 & 0.023 & & & 0.98 \\
$cay$                       & 726 & $-0.008$ & 0.039 & & & 0.99 \\[4pt]
\multicolumn{7}{l}{\textit{Panel C: Sample coverage (avg.\ stocks per month by decade)}} \\
1963--1969 & \multicolumn{2}{l}{1,425} & 1980--1989 & \multicolumn{2}{l}{2,750} & \\
1970--1979 & \multicolumn{2}{l}{2,205} & 1990--1999 & \multicolumn{2}{l}{3,358} & \\
2000--2009 & \multicolumn{2}{l}{2,645} & 2010--2023 & \multicolumn{2}{l}{2,168} & \\
\bottomrule
\end{tabular}
\begin{minipage}{\textwidth}
\vspace{4pt}
\footnotesize
\textit{Notes:} Panel~A reports the distribution of monthly excess returns (\%) and selected cross-sectionally ranked firm characteristics (rank $\in [0,1]$) across all stock-month observations in the sample (July 1963--December 2023). Panel~B reports time-series summary statistics for macroeconomic state variables at monthly frequency; AC(1) is the first-order autocorrelation. VIX is available from 1990; HKM intermediary capital ratio from 1970; Baker-Wurgler sentiment from 1965. Panel~C reports the average number of stocks per month within each decade.
\end{minipage}
\end{table}

% ══════════════════════════════════════════════════════════════════════
%  Table 2: Consumption-Based Euler Equation Restrictions (from Section 3)
% ══════════════════════════════════════════════════════════════════════

\begin{table}[h!]
\centering
\caption{Consumption-Based Euler Equation Restrictions}
\label{tab:euler}
\footnotesize\setlength{\tabcolsep}{3pt}
\begin{tabularx}{\textwidth}{c >{\raggedright\arraybackslash}p{3cm} >{\raggedright\arraybackslash}p{3cm} X}
\toprule
\# & Restriction & Source & SDF / Penalty Form \\
\midrule
1 & Power utility (C-CAPM)
  & \citep{breeden1979intertemporal,lucas1978asset}
  & $M_{t+1} = \delta\left(\frac{C_{t+1}}{C_t}\right)^{-\gamma}$ \\[6pt]
2 & Epstein--Zin recursive utility
  & \citep{epsteinzin1989substitution}
  & $M_{t+1} = \delta^\theta \left(\frac{C_{t+1}}{C_t}\right)^{-\theta/\psi} R_{m,t+1}^{\theta-1}$, \quad $\theta = \frac{1-\gamma}{1-1/\psi}$ \\[6pt]
3 & External habit (surplus ratio)
  & \citep{campbellcochrane1999}
  & $M_{t+1} = \delta\left(\frac{S_{t+1}}{S_t}\right)^{-\gamma}\left(\frac{C_{t+1}}{C_t}\right)^{-\gamma}$, \quad $S_t = \frac{C_t - H_t}{C_t}$ \\[6pt]
4 & Internal habit (difference)
  & \citep{constantinides1990habit}
  & $M_{t+1} = \delta\,\frac{(C_{t+1} - bC_t)^{-\gamma}}{(C_t - bC_{t-1})^{-\gamma}}$ \\[6pt]
5 & Internal habit (ratio)
  & \citep{abel1990asset}
  & $M_{t+1} = \delta\,\frac{(C_{t+1}/C_t^b)^{-\gamma}}{(C_t/C_{t-1}^b)^{-\gamma}}$ \\[6pt]
6 & Long-run risk
  & \citep{bansalyaron2004}
  & $M_{t+1}$ under Epstein--Zin with $\Delta c_{t+1} = \mu + x_t + \sigma_t \eta_{t+1}$; $x_t$ persistent, $\sigma_t^2$ stochastic \\[6pt]
7 & Long-run risk with jumps
  & \citep{drechsleryaron2011}
  & Adds compound Poisson jumps in $\sigma_t^2$ to the Bansal--Yaron SDF \\[6pt]
8 & Rare disaster (constant intensity)
  & \citep{barro2006rare,gabaix2012variable}
  & $M_{t+1}$ includes disaster probability $p$ and contraction size $b$: $M_{t+1} \propto (1 - p + p \cdot b^{-\gamma})$ \\[6pt]
9 & Variable rare disaster
  & \citep{wachter2013disaster}
  & Time-varying disaster intensity $\lambda_t$ following a Cox process; $M_{t+1}$ depends on $\lambda_t, \lambda_{t+1}$ \\[6pt]
10 & Precautionary savings (prudence)
  & \citep{kimball1990precautionary,leland1968saving}
  & $M_{t+1} \approx \delta(C_{t+1}/C_t)^{-\gamma}[1 + \frac{\gamma+1}{2}\sigma^2_{c,t}]$; prudence terms \\[6pt]
11 & Durable consumption
  & \citep{yogo2006durable}
  & Epstein--Zin with nondurable and durable consumption; $M_{t+1}$ depends on $(C_{t+1}/C_t)$, $(D_{t+1}/D_t)$, and $R_{m,t+1}$ \\[6pt]
12 & Heterogeneous agents
  & \citep{constantinidesduffie1996}
  & $M_{t+1} = \delta(C_{t+1}/C_t)^{-\gamma} \exp[-\frac{\gamma(\gamma+1)}{2}\sigma^2_{\text{cross}}]$; $\sigma^2_{\text{cross}}$ = cross-sectional cons.\ variance \\[6pt]
13 & Nonparametric Euler
  & \citep{escanciano2016nonparametric}
  & $M_{t+1} = g(C_{t+1}/C_t)$ with $g$ unknown, identified nonparametrically from Euler equation \\
\bottomrule
\end{tabularx}
\end{table}

% ══════════════════════════════════════════════════════════════════════
%  Table 3: Production-Based and Investment Restrictions (from Section 3)
% ══════════════════════════════════════════════════════════════════════

\begin{table}[h!]
\centering
\caption{Production-Based and Investment Restrictions}
\label{tab:production}
\footnotesize\setlength{\tabcolsep}{3pt}
\begin{tabularx}{\textwidth}{c >{\raggedright\arraybackslash}p{2.5cm} >{\raggedright\arraybackslash}p{2.5cm} X}
\toprule
\# & Restriction & Source & Explicit Penalty Form \\
\midrule
14 & Investment CAPM
   & \citep{cochrane1991production,liuwhitedzhang2009}
   & $C_{14}(f) = \frac{1}{n}\sum_s \bigl(f(\mathbf{X}_s) - \hat{a}_0 - \hat{a}_1 (I/K)_s - \hat{a}_2\,\mathrm{ROE}_s\bigr)^2$ \\[6pt]
15 & $q$-theory factor model
   & \citep{houxuezhang2015}
   & $C_{15}(f) = \frac{1}{n}\sum_s \bigl(f(\mathbf{X}_s) - \hat{a}_0 - \hat{a}_1 (I/K)_s - \hat{a}_2\,\mathrm{ROE}_s - \hat{a}_3\,\mathrm{EG}_s\bigr)^2$ \\[6pt]
16 & Convex adjustment costs
   & \citep{hayashi1982tobin}
   & $C_{16}(f) = \frac{1}{n}\sum_s \bigl(f(\mathbf{X}_s) - \hat{\phi}_1 (I/K)_s - \hat{\phi}_2 (I/K)_s^2\bigr)^2$ \\[6pt]
17 & Nonconvex adjustment costs
   & \citep{abeleberly1996,cooperhaltiwanger2006}
   & $C_{17}(f) = \frac{1}{n}\sum_s \bigl(f(\mathbf{X}_s) - \hat{a}_1 (I/K)_s - \hat{a}_2\,\ind_{\{|I/K|>h\}}\bigr)^2$ \\[6pt]
18 & Irreversibility
   & \citep{abeleberly1996,zhang2005value}
   & $C_{18}(f) = \frac{1}{n}\sum_s \bigl(f(\mathbf{X}_s) - \hat{a}_1 (I/K)_s - \hat{a}_2\,\ind_{\{I<0\}}\bigr)^2$ \\[6pt]
19 & Financial constraints (collateral)
   & \citep{kiyotakimoore1997}
   & $C_{19}(f) = \frac{1}{n}\sum_s \bigl(f(\mathbf{X}_s) - \hat{a}_0 - \hat{a}_1 (I/K)_s - \hat{a}_2 (K/\text{debt})_s\bigr)^2$ \\[6pt]
20 & Financial constraints (equity)
   & \citep{hennessywhited2005}
   & $C_{20}(f) = \frac{1}{n}\sum_s \bigl(f(\mathbf{X}_s) - \hat{a}_0 - \hat{a}_1 (I/K)_s - \hat{a}_2\,\mathrm{Lev}_s - \hat{a}_3\,\mathrm{CF}_s\bigr)^2$ \\[6pt]
21 & Costly external finance
   & \citep{livdansaprizazhang2009}
   & $C_{21}(f) = \frac{1}{n}\sum_s \bigl(f(\mathbf{X}_s) - \hat{a}_0 - \hat{a}_1 (I/K)_s - \hat{a}_2\,\mathrm{Lev}_s - \hat{a}_3\,\mathrm{Div}_s\bigr)^2$ \\[6pt]
22 & Labor-augmented
   & \citep{belolinbazdresch2014}
   & $C_{22}(f) = \frac{1}{n}\sum_s \bigl(f(\mathbf{X}_s) - \hat{a}_0 - \hat{a}_1 (I/K)_s - \hat{a}_2\,\mathrm{HN}_s\bigr)^2$ \\[6pt]
23 & Intangible capital
   & \citep{eisfeldtpapanikolaou2013}
   & $C_{23}(f) = \frac{1}{n}\sum_s \bigl(f(\mathbf{X}_s) - \hat{a}_0 - \hat{a}_1 (I/K)_s - \hat{a}_2 (\text{R\&D}/K)_s - \hat{a}_3 (\text{SGA}/K)_s\bigr)^2$ \\
\bottomrule
\end{tabularx}
\end{table}

% ══════════════════════════════════════════════════════════════════════
%  Table 4: Intermediary and Institutional Restrictions (from Section 3)
% ══════════════════════════════════════════════════════════════════════

\begin{table}[h!]
\centering
\caption{Intermediary and Institutional Restrictions}
\label{tab:intermediary}
\footnotesize\setlength{\tabcolsep}{3pt}
\begin{tabularx}{\textwidth}{c >{\raggedright\arraybackslash}p{2.5cm} >{\raggedright\arraybackslash}p{2.5cm} X}
\toprule
\# & Restriction & Source & SDF / Penalty Form \\
\midrule
24 & Intermediary capital ratio
   & \citep{hekrishnamurthy2013}
   & $M_{24,t+1} = \hat{a}_{24} + \hat{b}_{24}\,\eta_{t+1}$; \quad $C_{24}(f) = \sum_t \bigl(\frac{1}{N}\sum_i M_{24,t+1}(f(\mathbf{X}_{i,t}) + R_t^f) - 1\bigr)^2$ \\[6pt]
25 & Intermediary leverage
   & \citep{adrianetulamuir2014}
   & $M_{25,t+1} = \hat{a}_{25} + \hat{b}_{25}\,\Delta\mathrm{Lev}_{t+1}$; same Euler penalty form \\[6pt]
26 & Funding liquidity
   & \citep{brunnermeier2009}
   & $M_{26,t+1} = \hat{a}_{26} + \hat{b}_{26}\,\mathrm{TED}_{t+1}$; margin spirals amplify funding shocks \\[6pt]
27 & Dealer balance sheet
   & \citep{hekellymanela2017}
   & $M_{27,t+1} = \hat{a}_{27} + \hat{b}_{27}\,\eta^{\mathrm{PD}}_{t+1}$; primary dealer equity capital ratio \\[6pt]
28 & Limits to arbitrage
   & \citep{delong1990noise}
   & $M_{28,t+1} = \hat{a}_{28} + \hat{b}_{28}\,\mathrm{Sent}_{t+1}$; noise-trader sentiment as SDF factor \\[6pt]
29 & Slow-moving capital
   & \citep{duffie2010slow}
   & $M_{29,t+1} = \hat{a}_{29} + \hat{b}_{29}\,\mathrm{Flow}_{t+1}$; capital flow proxy \\[6pt]
30 & Margin CAPM
   & \citep{garleanupedersen2011}
   & $M_{30,t+1} = \hat{a}_{30} + \hat{b}_{30}^m R_{m,t+1} + \hat{b}_{30}^{\psi}\,\psi_{t+1}$; margin requirement factor \\[6pt]
31 & Leverage constraints
   & \citep{frazzinipedersen2014}
   & $M_{31,t+1} = \hat{a}_{31} + \hat{b}_{31}\,\mathrm{BAB}_{t+1}$; betting-against-beta factor \\
\bottomrule
\end{tabularx}
\end{table}

% ══════════════════════════════════════════════════════════════════════
%  Table 5: Information and Learning Restrictions (from Section 3)
% ══════════════════════════════════════════════════════════════════════

\begin{table}[h!]
\centering
\caption{Information and Learning Restrictions}
\label{tab:information}
\footnotesize\setlength{\tabcolsep}{3pt}
\begin{tabularx}{\textwidth}{c >{\raggedright\arraybackslash}p{2.5cm} >{\raggedright\arraybackslash}p{2.5cm} X}
\toprule
\# & Restriction & Source & SDF / Penalty Form \\
\midrule
32 & Bayesian learning about growth
   & \citep{pastorveronesi2009}
   & $M_{32,t+1}$ depends on posterior mean $\hat{\mu}_t$ and variance $\hat{\sigma}^2_t$ of long-run growth; $C_{32}(f) = \sum_t (\frac{1}{N}\sum_i M_{32,t+1}(f(\mathbf{X}_{i,t}) + R_t^f) - 1)^2$ \\[6pt]
33 & Learning about regime
   & \citep{david2008heterogeneous,veronesi1999}
   & $M_{33,t+1}$ depends on regime probability $\pi_t = P(\text{good} \mid \Fcal_t)$; belief updates amplify volatility \\[6pt]
34 & Parameter uncertainty
   & \citep{johanneskorteweg2014}
   & $M_{34,t+1} = \int M(\theta) \, dP(\theta \mid \Fcal_t)$; predictive variance enters pricing \\[6pt]
35 & Ambiguity aversion
   & \citep{epsteinschneider2008,ilutschneider2014}
   & $M_{35,t+1} = \min_{P \in \mathcal{P}} dP/dP_0$; worst-case density ratio as SDF \\[6pt]
36 & Disagreement
   & \citep{basak2005asset,xiongyan2010}
   & $M_{36,t+1} = \hat{a}_{36} + \hat{b}_{36}\,\mathrm{Disp}_{t+1}$; belief dispersion as linear SDF factor \\[6pt]
37 & Rational inattention
   & \citep{mackowiak2009optimal,vannieuwerburgh2010information}
   & $M_{37,t+1} = \hat{a}_{37} + \hat{b}_{37}\,\mathrm{InfoFlow}_{t+1}$; information capacity as linear SDF factor \\
\bottomrule
\end{tabularx}
\end{table}

% ══════════════════════════════════════════════════════════════════════
%  Table 6: Demand-System and Market Microstructure Restrictions (from Section 3)
% ══════════════════════════════════════════════════════════════════════

\begin{table}[h!]
\centering
\caption{Demand-System and Market Microstructure Restrictions}
\label{tab:demand}
\footnotesize\setlength{\tabcolsep}{3pt}
\begin{tabularx}{\textwidth}{c >{\raggedright\arraybackslash}p{2.5cm} >{\raggedright\arraybackslash}p{2.5cm} X}
\toprule
\# & Restriction & Source & SDF / Penalty Form \\
\midrule
38 & Market clearing
   & \citep{koijenyogo2019}
   & $C_{38}(f) = \sum_t \|\sum_i w_i D_i(\hat{\mathbf{f}}_t) - \mathbf{S}_t\|^2$; aggregate demand = supply \\[6pt]
39 & Portfolio choice FOC
   & \citep{koijenrichmondyogo2024}
   & $C_{39}(f) = \sum_t \|\sum_i w_i D_i^{\mathrm{FOC}}(\hat{\mathbf{f}}_t, \mathbf{X}_t) - \mathbf{S}_t\|^2$; characteristic-based elasticities \\[6pt]
40 & Inelastic markets
   & \citep{gabaixkoijen2021}
   & $C_{40}(f) = \sum_t \|\hat{\epsilon}^{-1}\,\Delta\mathrm{Flow}_t - \Delta\hat{\mathbf{f}}_t\|^2$; flow-driven price impact \\[6pt]
41 & Search frictions OTC
   & \citep{duffiegarleanupedersen2005}
   & $C_{41}(f) = \frac{1}{n}\sum_s (f(\mathbf{X}_s) - \hat{a}_0 - \hat{a}_1\,\lambda_s - \hat{a}_2\,\beta^{\mathrm{barg}}_s)^2$ \\[6pt]
42 & Fire sale externalities
   & \citep{shleifervishny1992}
   & $M_{42,t+1} = \hat{a}_{42} + \hat{b}_{42}\,\mathrm{FSP}_{t+1}$; fire-sale pressure as linear SDF \\[6pt]
43 & Short-sale constraints
   & \citep{miller1977risk,hongstein2003}
   & $C_{43}(f) = \frac{1}{n}\sum_s (f(\mathbf{X}_s) - \hat{a}_0 - \hat{a}_1\,\mathrm{SI}_s - \hat{a}_2\,\mathrm{Disp}_s)^2$ \\
\bottomrule
\end{tabularx}
\end{table}

% ══════════════════════════════════════════════════════════════════════
%  Table 7: Variance and Volatility Restrictions (from Section 3)
% ══════════════════════════════════════════════════════════════════════

\begin{table}[h!]
\centering
\caption{Variance and Volatility Restrictions}
\label{tab:options}
\footnotesize\setlength{\tabcolsep}{3pt}
\begin{tabularx}{\textwidth}{c >{\raggedright\arraybackslash}p{2.5cm} >{\raggedright\arraybackslash}p{2.5cm} X}
\toprule
\# & Restriction & Source & Explicit Penalty Form \\
\midrule
44 & Variance risk premium
   & \citep{bollerslevtauchenzhou2009}
   & $C_{44}(f) = \frac{1}{n}\sum_s (f(\mathbf{X}_s) - \hat{a}_0 - \hat{a}_1\,\mathrm{VRP}_{t(s)})^2$; $\mathrm{VRP}_t = \mathrm{IV}_t^2 - \hat{\E}_t[\mathrm{RV}_{t+1}]$ \\[6pt]
45 & Stochastic volatility
   & \citep{heston1993}
   & $C_{45}(f) = \frac{1}{n}\sum_s (f(\mathbf{X}_s) - \hat{a}_0 - \hat{a}_1\,V_{t(s)})^2$; $V_t$ = conditional variance proxy (VIX$^2$) \\
\bottomrule
\end{tabularx}
\end{table}

% ══════════════════════════════════════════════════════════════════════
%  Table 8: Behavioral and Sentiment Restrictions (from Section 3)
% ══════════════════════════════════════════════════════════════════════

\begin{table}[h!]
\centering
\caption{Behavioral and Sentiment Restrictions}
\label{tab:behavioral}
\footnotesize\setlength{\tabcolsep}{3pt}
\begin{tabularx}{\textwidth}{c >{\raggedright\arraybackslash}p{2.5cm} >{\raggedright\arraybackslash}p{2.5cm} X}
\toprule
\# & Restriction & Source & SDF / Penalty Form \\
\midrule
46 & Prospect theory
   & \citep{barberishuangsantos2001}
   & $M_{46,t+1} = \delta(C_{t+1}/C_t)^{-\gamma}[1 + \hat{\lambda}\,\max(0, -R_{m,t+1})/R_{m,t}]$; loss aversion kink \\[6pt]
47 & Overconfidence
   & \citep{danielhirshleifer1998}
   & $C_{47}(f) = \frac{1}{n}\sum_s (f(\mathbf{X}_s) - \hat{a}_0 - \hat{a}_1\,\mathrm{Mom}_s - \hat{a}_2\,\mathrm{Rev}_s)^2$ \\[6pt]
48 & Disposition effect
   & \citep{frazzini2006disposition}
   & $C_{48}(f) = \frac{1}{n}\sum_s (f(\mathbf{X}_s) - \hat{a}_0 - \hat{a}_1\,\mathrm{CGO}_s)^2$ \\[6pt]
49 & Extrapolative expectations
   & \citep{barberis2018extrapolation}
   & $C_{49}(f) = \frac{1}{n}\sum_s (f(\mathbf{X}_s) - \hat{a}_0 - \hat{a}_1\,\bar{R}^{\text{past}}_s)^2$ \\[6pt]
50 & Investor sentiment
   & \citep{bakerwurgler2006}
   & $M_{50,t+1} = \hat{a}_{50} + \hat{b}_{50}\,\mathrm{Sent}_{t+1}$; high sentiment predicts low returns for speculative stocks \\[6pt]
51 & Salience theory
   & \citep{bordalo2012salience}
   & $C_{51}(f) = \frac{1}{n}\sum_s (f(\mathbf{X}_s) - \hat{a}_0 - \hat{a}_1\,\mathrm{Sal}_s)^2$ \\
\bottomrule
\end{tabularx}
\end{table}

% ══════════════════════════════════════════════════════════════════════
%  Table 9: Macro-Finance and Term Structure Restrictions (from Section 3)
% ══════════════════════════════════════════════════════════════════════

\begin{table}[h!]
\centering
\caption{Macro-Finance and Term Structure Restrictions}
\label{tab:macro}
\footnotesize\setlength{\tabcolsep}{3pt}
\begin{tabularx}{\textwidth}{c >{\raggedright\arraybackslash}p{2.5cm} >{\raggedright\arraybackslash}p{2.5cm} X}
\toprule
\# & Restriction & Source & SDF / Penalty Form \\
\midrule
52 & Affine term structure
   & \citep{duffiekan1996,angpiazzesi2003}
   & $C_{52}(f) = \frac{1}{n}\sum_s (f(\mathbf{X}_s) - \hat{a}_0 - \hat{\mathbf{a}}^\top \mathbf{y}_{t(s)})^2$; yield curve state $\mathbf{y}_t$ \\[6pt]
53 & Inflation risk premium
   & \citep{bekaertengstrom2010}
   & $M_{53,t+1} = \hat{a}_{53} + \hat{b}_{53}\,\mathrm{IRP}_{t+1}$; inflation risk premium as linear SDF factor \\[6pt]
54 & Monetary policy
   & \citep{bernankekuttner2005}
   & $C_{54}(f) = \frac{1}{n}\sum_s (f(\mathbf{X}_s) - \hat{a}_0 - \hat{a}_1\,\Delta i^{\text{surp}}_{t(s)})^2$ \\[6pt]
55 & Consumption-wealth ratio
   & \citep{lettauludvigson2001}
   & $C_{55}(f) = \frac{1}{n}\sum_s (f(\mathbf{X}_s) - \hat{a}_0 - \hat{a}_1\,cay_{t(s)})^2$ \\[6pt]
56 & Credit spread / EBP
   & \citep{gilchristzakrajsek2012}
   & $M_{56,t+1} = \hat{a}_{56} + \hat{b}_{56}\,\mathrm{EBP}_{t+1}$; excess bond premium as linear SDF factor \\
\bottomrule
\end{tabularx}
\end{table}

% ══════════════════════════════════════════════════════════════════════
%  Table 10: Out-of-Sample Prediction Performance (from Section 6)
% ══════════════════════════════════════════════════════════════════════

\begin{table}[h!]
\centering
\caption{Out-of-Sample Prediction Performance (Managed Portfolios)}
\label{tab:oos_performance}
\small
\begin{tabular}{lcc}
\toprule
Model & $R^2_{\text{OOS}}$ (\%) & DM vs.\ Hist.\ Mean \\
\midrule
\multicolumn{3}{l}{\textit{Linear benchmarks}} \\
Historical mean           & $-0.14$ & ---     \\
OLS (all predictors)      & $32.05$ & $18.91$ \\
Ridge regression          & $32.05$ & $18.91$ \\
Lasso                     & $31.66$ & $19.00$ \\
Elastic net               & $31.86$ & $18.97$ \\[4pt]
\multicolumn{3}{l}{\textit{Nonlinear benchmarks}} \\
Kernel ridge regression   & $27.86$ & $19.66$ \\
Random forest             & $36.75$ & $18.85$ \\[4pt]
\multicolumn{3}{l}{\textit{Theory-informed models}} \\
Theory-KRR (all groups)   & $27.86$ & $19.66$ \\
\bottomrule
\end{tabular}
\begin{minipage}{\textwidth}
\vspace{4pt}
\footnotesize
\textit{Notes:} $R^2_{\text{OOS}}$ is defined as $1 - \sum_s (R_s - \hat{f}(\mathbf{X}_s))^2 / \sum_s (R_s - \bar{R})^2$, where $\bar{R}$ is the historical mean computed using only data available at the time of prediction. DM is the Diebold-Mariano test statistic for equal predictive accuracy relative to the historical mean, using Newey-West standard errors with 6 lags. The evaluation period is 1983--2023. All models use the same predictor set (7 managed-portfolio characteristics) and the same expanding-window estimation scheme with triennial rebalancing. These are preliminary results with fixed hyperparameters; cross-validated results with individually tuned $\bm{\mu}$ are forthcoming.
\end{minipage}
\end{table}

% ══════════════════════════════════════════════════════════════════════
%  Table 11: Temporal Variation in Optimal Theory Family (from Section 6)
% ══════════════════════════════════════════════════════════════════════

\begin{table}[h!]
\centering
\caption{Temporal Variation in Optimal Theory Family}
\label{tab:theory_ranking}
\small
\begin{tabular}{llll}
\toprule
OOS Period & Best Theory Family & Top Restrictions & Val.\ MSE gain (\%) \\
\midrule
1983--1985  & All groups ($\mu = 0.1$)   & Realized vol, demand elasticity  & 0.030 \\
1986--1988  & Demand systems             & Demand elasticity, substitution  & 0.017 \\
1989--1991  & Macro-finance              & Term spread, default spread, EBP & 0.000 \\
1992--1994  & Volatility                 & Realized volatility              & 0.066 \\
1995--1997  & Consumption + Intermediary & Precautionary savings, $cay$, HKM capital  & 0.091 \\
1998--2000  & Production                 & Investment, leverage, asset growth & 0.022 \\
2001--2003  & Behavioral                 & Momentum, LT reversal, disposition & 0.041 \\
2004--2006  & All groups ($\mu = 1.0$)   & Realized vol, precautionary savings & 0.005 \\
2007--2009  & No production ($\mu = 0.1$) & Realized vol, $cay$, demand      & 0.016 \\
2010--2012  & Macro-finance              & Term spread, default spread, EBP & 0.024 \\
\bottomrule
\end{tabular}
\begin{minipage}{\textwidth}
\vspace{4pt}
\footnotesize
\textit{Notes:} Each row reports the best-performing $\mu$ configuration for the indicated OOS period, selected by temporal validation on the last 24~months of the training set. Val.\ MSE gain is the percentage reduction in validation MSE relative to standard KRR ($\mu_j = 0$ for all $j$). The evaluation uses an expanding window with 240-month minimum training period and triennial rebalancing.
\end{minipage}
\end{table}

% ══════════════════════════════════════════════════════════════════════
%  Table 12: Theory-KRR Validation Improvement by Subperiod (from Section 6)
% ══════════════════════════════════════════════════════════════════════

\begin{table}[h!]
\centering
\caption{Theory-KRR Validation Improvement over KRR by Subperiod}
\label{tab:subsample}
\small
\begin{tabular}{lccc}
\toprule
Subperiod & Windows & Avg.\ Val.\ MSE Gain (\%) & Best Theory Family \\
\midrule
Pre-2000 (1983--2000)  & 0--5  & 0.037 & Demand, Vol, Cons+Inter, Prod \\
Post-2000 (2001--2012) & 6--9  & 0.022 & Behavioral, All, No-prod, Macro \\[4pt]
Crisis windows         & 6, 8  & 0.029 & Behavioral (dot-com), No-prod (GFC) \\
Non-crisis windows     & 0--5, 7, 9 & 0.029 & Varies by period \\[4pt]
Early training ($<$300mo) & 0--1 & 0.024 & All groups, Demand \\
Late training ($>$400mo)  & 5--9 & 0.022 & Prod, Behav, All, No-prod, Macro \\
\bottomrule
\end{tabular}
\begin{minipage}{\textwidth}
\vspace{4pt}
\footnotesize
\textit{Notes:} Val.\ MSE gain is the percentage reduction in validation MSE of the best Theory-KRR configuration relative to standard KRR. Crisis windows are those whose OOS period includes the dot-com bust (2001--2003) or the global financial crisis (2007--2009). Improvement is measured on the temporal validation set within each window before out-of-sample prediction.
\end{minipage}
\end{table}


% ══════════════════════════════════════════════════════════════════════
%  Appendix
% ══════════════════════════════════════════════════════════════════════

\newpage
\appendix
%!TEX root = ../main.tex
%----------------------------------------------------------------------
%  Appendix
%  Paper:   Theory-Informed Kernel Ridge Regression for Asset Pricing
%  Authors: Amedeo Andriollo and Juan F. Imbet
%----------------------------------------------------------------------

\section{Proofs}
\label{app:proofs}

%----------------------------------------------------------------------
\subsection{Proof of Proposition~\ref{prop:representer} (Representer Theorem)}
%----------------------------------------------------------------------

\begin{proof}
Decompose any $f \in \mathcal{H}$ as $f = f_{\parallel} + f_{\perp}$,
where $f_{\parallel}$ lies in the finite-dimensional subspace
$\mathcal{V} = \operatorname{span}\{k(\mathbf{X}_s, \cdot) : s = 1, \ldots, n\}$
and $f_{\perp} \in \mathcal{V}^{\perp}$.
By the reproducing property,
\[
f(\mathbf{X}_s)
= \langle f, k(\mathbf{X}_s, \cdot) \rangle_{\mathcal{H}}
= \langle f_{\parallel}, k(\mathbf{X}_s, \cdot) \rangle_{\mathcal{H}}
  + \langle f_{\perp}, k(\mathbf{X}_s, \cdot) \rangle_{\mathcal{H}}
= f_{\parallel}(\mathbf{X}_s),
\]
since $k(\mathbf{X}_s, \cdot) \in \mathcal{V}$ and
$f_{\perp} \perp \mathcal{V}$.
Therefore the prediction loss
$\frac{1}{n}\sum_s (R_s - f(\mathbf{X}_s))^2$
depends only on $f_{\parallel}$, and so does every $C_j(f)$
by assumption~\eqref{eq:penalty_eval}.

For the RKHS norm, orthogonality gives
$\|f\|_{\mathcal{H}}^2 = \|f_{\parallel}\|_{\mathcal{H}}^2 +
\|f_{\perp}\|_{\mathcal{H}}^2$.
Since $\|f_{\perp}\|_{\mathcal{H}}^2 \geq 0$ and
$\lambda_{\mathrm{stat}} > 0$, the objective is strictly decreased
by setting $f_{\perp} = 0$ whenever $f_{\perp} \neq 0$.
Hence any minimizer satisfies $f = f_{\parallel} \in \mathcal{V}$,
which yields the representation~\eqref{eq:representer}.
\end{proof}

%----------------------------------------------------------------------
\subsection{Proof of Proposition~\ref{prop:finite_dim} (Finite-Dimensional Problem)}
%----------------------------------------------------------------------

\begin{proof}
The prediction loss becomes
$\frac{1}{n}\|\mathbf{R} - \mathbf{K}\boldsymbol{\alpha}\|^2$
since $f(\mathbf{X}_s) = \sum_{s'} \alpha_{s'} k(\mathbf{X}_{s'}, \mathbf{X}_s)
= [\mathbf{K}\boldsymbol{\alpha}]_s$.
For the RKHS norm,
\[
\|f\|_{\mathcal{H}}^2
= \left\langle \sum_s \alpha_s\, k(\mathbf{X}_s, \cdot),\;
         \sum_{s'} \alpha_{s'}\, k(\mathbf{X}_{s'}, \cdot) \right\rangle_{\mathcal{H}}
= \sum_{s,s'} \alpha_s\, \alpha_{s'}\, k(\mathbf{X}_s, \mathbf{X}_{s'})
= \boldsymbol{\alpha}^\top \mathbf{K}\, \boldsymbol{\alpha}.
\]
By condition~\eqref{eq:penalty_eval}, each $C_j(f)$ is a function of
$\hat{\mathbf{f}} = \mathbf{K}\boldsymbol{\alpha}$.
Combining yields~\eqref{eq:finite_dim}.
\end{proof}

%----------------------------------------------------------------------
\subsection{Proof of Proposition~\ref{prop:foc} (First-Order Conditions)}
%----------------------------------------------------------------------

\begin{proof}
Write the objective~\eqref{eq:finite_dim} with quadratic penalties as
\begin{multline}
\label{eq:obj_expanded}
L(\boldsymbol{\alpha})
= \frac{1}{n}(\mathbf{R} - \mathbf{K}\boldsymbol{\alpha})^\top
  (\mathbf{R} - \mathbf{K}\boldsymbol{\alpha})
+ \lambda_{\mathrm{stat}}\,\boldsymbol{\alpha}^\top\mathbf{K}\boldsymbol{\alpha} \\
+ \sum_{j \in \mathcal{J}^A} \mu_j
  (\mathbf{A}_j\mathbf{K}\boldsymbol{\alpha} - \mathbf{b}_j)^\top
  (\mathbf{A}_j\mathbf{K}\boldsymbol{\alpha} - \mathbf{b}_j)
+ \sum_{j \in \mathcal{J}^{B}} \frac{\mu_j}{n}
  (\mathbf{K}\boldsymbol{\alpha} - \mathbf{v}_j)^\top
  (\mathbf{K}\boldsymbol{\alpha} - \mathbf{v}_j),
\end{multline}
where $\mathbf{v}_j = \mathbf{g}_j$ for Type~B.
Taking the gradient with respect to $\boldsymbol{\alpha}$:
\begin{align}
\frac{\partial L}{\partial \boldsymbol{\alpha}}
&= -\frac{2}{n}\,\mathbf{K}(\mathbf{R} - \mathbf{K}\boldsymbol{\alpha})
+ 2\lambda_{\mathrm{stat}}\,\mathbf{K}\boldsymbol{\alpha}
+ 2\sum_{j \in \mathcal{J}^A} \mu_j\,\mathbf{K}\mathbf{A}_j^\top
  (\mathbf{A}_j\mathbf{K}\boldsymbol{\alpha} - \mathbf{b}_j)
\notag \\
&\quad + 2\sum_{j \in \mathcal{J}^{B}} \frac{\mu_j}{n}\,\mathbf{K}
  (\mathbf{K}\boldsymbol{\alpha} - \mathbf{v}_j).
\label{eq:gradient}
\end{align}
Setting~\eqref{eq:gradient} to zero and left-multiplying by $\mathbf{K}^{-1}$
(which exists when $\mathbf{K} \succ 0$):
\[
\frac{1}{n}\,\mathbf{K}\boldsymbol{\alpha}
+ \lambda_{\mathrm{stat}}\,\boldsymbol{\alpha}
+ \sum_{j \in \mathcal{J}^A} \mu_j\,\mathbf{A}_j^\top\mathbf{A}_j
  \mathbf{K}\boldsymbol{\alpha}
+ \sum_{j \in \mathcal{J}^{B}} \frac{\mu_j}{n}\,\mathbf{K}\boldsymbol{\alpha}
= \frac{1}{n}\,\mathbf{R}
+ \sum_{j \in \mathcal{J}^A} \mu_j\,\mathbf{A}_j^\top\mathbf{b}_j
+ \sum_{j \in \mathcal{J}^{B}} \frac{\mu_j}{n}\,\mathbf{v}_j.
\]
Collecting terms using the definitions of $\boldsymbol{\Omega}$
and $\boldsymbol{\omega}$ from~\eqref{eq:Omega}--\eqref{eq:omega_vec}
gives~\eqref{eq:foc}.
Strict positive definiteness of $\mathbf{K}$ and $\lambda_{\mathrm{stat}} > 0$
ensure the coefficient matrix in~\eqref{eq:foc} is invertible, yielding
uniqueness.
\end{proof}

%----------------------------------------------------------------------
\subsection{Proof of Proposition~\ref{prop:nesting} (Nesting)}
%----------------------------------------------------------------------

\begin{proof}
Claims~(i) and~(ii) follow directly from the definition.
For~(i), setting $\mu_j = 0$ eliminates the structural penalties.
For~(ii), taking $\mu_j \to \infty$ enforces $C_j(f) = 0$ as a constraint;
formally, the sequence of penalized minimizers converges to the
constrained minimizer provided $C_j$ is lower semicontinuous and
$\{f \in \mathcal{H} : C_j(f) = 0\}$ is nonempty.
Setting $\lambda_{\mathrm{stat}} = 0$ removes agnostic regularization.
\end{proof}

%----------------------------------------------------------------------
\subsection{Proof of Proposition~\ref{prop:misspecification} (Misspecification Robustness)}
%----------------------------------------------------------------------

\begin{proof}
We show that any $\mu_j > 0$ introduces an asymptotic bias, so the
oracle sets $\mu_j^* = 0$.

\emph{Step 1: Asymptotic bias from misspecification.}
When $C_j(f^*) > 0$, the set
$\mathcal{F}_j(\mu_j) = \arg\min_{f \in \mathcal{H}}
\frac{1}{n}\sum_s (R_s - f(\mathbf{X}_s))^2
+ \lambda_{\mathrm{stat}}\|f\|_{\mathcal{H}}^2 + \mu_j C_j(f)$
converges (as $n \to \infty$ and $\lambda_{\mathrm{stat}} \to 0$ at an
appropriate rate) to the population minimizer
\[
f_j^*(\mu_j) = \arg\min_{f} \mathbb{E}[(R - f(\mathbf{X}))^2] + \mu_j C_j(f).
\]
By~\ref{cond:consistency}, $f_j^*(0) = f^*$.
For $\mu_j > 0$, the term $\mu_j C_j(f)$ tilts the solution toward
$\{f : C_j(f) \text{ is small}\}$.
Since $C_j(f^*) > 0$, this tilt is binding: $f_j^*(\mu_j) \neq f^*$
for all $\mu_j > 0$, and the prediction risk satisfies
\[
\mathbb{E}\bigl[(R - f_j^*(\mu_j)(\mathbf{X}))^2\bigr]
> \mathbb{E}\bigl[(R - f^*(\mathbf{X}))^2\bigr]
= \mathbb{E}[\varepsilon^2]
\qquad \text{for all } \mu_j > 0.
\]

\emph{Step 2: Risk gap is bounded away from zero.}
By continuity of $C_j$ (condition~\ref{cond:bounded}) and strict
convexity of the population prediction risk,
for any $\delta > 0$ there exists $c(\delta) > 0$ such that
\[
\mathrm{Risk}(\mu_j) - \mathrm{Risk}(0)
\;\geq\; c(\delta) > 0
\qquad \text{for all } \mu_j \geq \delta,
\]
where $\mathrm{Risk}(\mu_j) = \mathbb{E}[(R - f_j^*(\mu_j)(\mathbf{X}))^2]$.

\emph{Step 3: Oracle efficiency implies $\hat{\mu}_j \to 0$.}
By~\ref{cond:cv_oracle}, $\mathrm{CV}(\hat{\boldsymbol{\lambda}})$
tracks the oracle $\inf_{\boldsymbol{\lambda}} \mathrm{CV}(\boldsymbol{\lambda})$
up to $o_p(1)$.  Since the oracle selects $\mu_j^* = 0$
(the unique risk minimizer along the $\mu_j$ coordinate by Step~2),
and the risk gap is bounded below by $c(\delta)$ for $\mu_j \geq \delta$,
we have $\Pr(\hat{\mu}_j \geq \delta) \to 0$ for every $\delta > 0$.
\end{proof}

%----------------------------------------------------------------------
\subsection{Proof of Proposition~\ref{prop:test_mu} (Testing the Marginal Value of Economic Models)}
\label{app:proof_test_mu}
%----------------------------------------------------------------------

\begin{proof}
\emph{Under $H_0$.}
When $\mu_j^* = 0$, the oracle solution does not use model $j$.
Constraining $\mu_j = 0$ therefore imposes no binding restriction on the
population objective.
The difference
$\mathrm{CV}(\hat{\boldsymbol{\lambda}}_{-j}) -
\mathrm{CV}(\hat{\boldsymbol{\lambda}})$
arises solely from finite-sample noise in the cross-validated selection
of the remaining hyperparameters.
By condition~\ref{cond:test_rate} and the oracle efficiency
in~\ref{cond:cv_oracle}, this difference is
$O_p(n^{-1} r_n)$, so $T_j = n \cdot O_p(n^{-1} r_n) = O_p(r_n) = o_p(1)$.

\emph{Under $H_1$.}
When $\mu_j^* > 0$, constraining $\mu_j = 0$ removes a binding
restriction.
By condition~\ref{cond:test_smooth},
the population risk increases by at least
$c > 0$ (a positive constant depending on $\mu_j^*$ and $C_j$).
By uniform convergence of the cross-validated risk to the population risk,
$\mathrm{CV}(\hat{\boldsymbol{\lambda}}_{-j}) -
\mathrm{CV}(\hat{\boldsymbol{\lambda}}) \geq c/2 > 0$
with probability approaching one.
Hence $T_j = n(c/2 + o_p(1)) \to \infty$.
\end{proof}


%======================================================================
\section{Computational Implementation}
\label{app:computation}
%======================================================================

This appendix provides implementation details for reproducing the estimation results.
All computations are performed on a workstation with a 40-core CPU, 274~GB of RAM, and an NVIDIA RTX A6000 GPU with 48~GB of VRAM.

%----------------------------------------------------------------------
\subsection{Nystr\"{o}m Approximation}
\label{app:nystrom}
%----------------------------------------------------------------------

The full kernel matrix $\mathbf{K} \in \mathbb{R}^{n \times n}$ requires $O(n^2)$ storage and $O(n^3)$ time for the Cholesky factorization.
For the managed-portfolio specification ($n \approx 4{,}000$), the full kernel is feasible.
For the individual-stock panel ($n > 100{,}000$), the full kernel matrix exceeds 80~GB and is infeasible even on large-memory workstations.

We address this via the Nystr\"{o}m low-rank approximation.
Select $m \ll n$ landmark points uniformly at random from the training data and compute:
\begin{align}
\mathbf{K}_{nm} &\in \mathbb{R}^{n \times m}, \qquad [\mathbf{K}_{nm}]_{ij} = k(\mathbf{X}_i, \mathbf{X}_{l_j}), \label{eq:nystrom_knm} \\
\mathbf{K}_{mm} &\in \mathbb{R}^{m \times m}, \qquad [\mathbf{K}_{mm}]_{ij} = k(\mathbf{X}_{l_i}, \mathbf{X}_{l_j}), \label{eq:nystrom_kmm}
\end{align}
where $\{l_1, \ldots, l_m\}$ are the landmark indices.
The Nystr\"{o}m feature map is $\mathbf{Z} = \mathbf{K}_{nm}\,\mathbf{K}_{mm}^{-1/2} \in \mathbb{R}^{n \times m}$, where $\mathbf{K}_{mm}^{-1/2}$ is obtained via eigendecomposition.
The kernel matrix is then approximated as $\mathbf{K} \approx \mathbf{Z}\,\mathbf{Z}^\top$.

Under this approximation, the KRR solution becomes
\begin{equation}
\label{eq:nystrom_solve}
\hat{\boldsymbol{\beta}} = (\mathbf{Z}^\top\mathbf{Z} + n\lambda_{\mathrm{stat}}\,\mathbf{I}_m)^{-1}\,\mathbf{Z}^\top\mathbf{y},
\end{equation}
which requires only an $m \times m$ Cholesky factorization at cost $O(m^3)$, plus $O(nm)$ for the matrix-vector products.
Predictions on new data $\mathbf{X}_{\mathrm{new}}$ are computed as $\hat{\mathbf{f}} = \mathbf{Z}_{\mathrm{new}}\,\hat{\boldsymbol{\beta}}$, where $\mathbf{Z}_{\mathrm{new}} = \mathbf{K}_{\mathrm{new},m}\,\mathbf{K}_{mm}^{-1/2}$.

For Theory-KRR with structural penalties, the L-BFGS optimization operates on the $m$-dimensional coefficient vector $\boldsymbol{\beta}$ rather than the $n$-dimensional $\boldsymbol{\alpha}$.
Each L-BFGS iteration computes the fitted values $\hat{\mathbf{f}} = \mathbf{Z}\boldsymbol{\beta}$ at cost $O(nm)$, evaluates the structural penalties $C_j(\hat{\mathbf{f}})$, and updates $\boldsymbol{\beta}$.
The gradient with respect to $\boldsymbol{\beta}$ has a simple form:
\begin{equation}
\label{eq:nystrom_gradient}
\frac{\partial L}{\partial \boldsymbol{\beta}} = -2\,\mathbf{Z}^\top(\mathbf{y} - \mathbf{Z}\boldsymbol{\beta}) + 2n\lambda_{\mathrm{stat}}\,\boldsymbol{\beta} + \sum_{j} \mu_j\,\mathbf{Z}^\top\frac{\partial C_j}{\partial \hat{\mathbf{f}}},
\end{equation}
where $\partial C_j / \partial \hat{\mathbf{f}}$ is the $n$-dimensional gradient of the $j$-th penalty with respect to the fitted values.

We use $m = 500$ landmarks throughout.
Table~\ref{tab:nystrom_scaling} summarizes the computational scaling.

\begin{table}[h]
\centering
\caption{Computational Scaling: Full Kernel vs.\ Nystr\"{o}m Approximation}
\label{tab:nystrom_scaling}
\small
\begin{tabular}{lcc}
\toprule
 & Full Kernel & Nystr\"{o}m ($m = 500$) \\
\midrule
Kernel storage                & $O(n^2)$  & $O(nm)$ \\
KRR solve                     & $O(n^3)$  & $O(nm^2 + m^3)$ \\
L-BFGS per iteration          & $O(n^2)$  & $O(nm)$ \\
\midrule
$n = 4{,}000$ (portfolios)    & 128~MB, $<1$s & 16~MB, $<0.1$s \\
$n = 90{,}000$ (stocks)       & 65~GB, $\sim 10$h & 360~MB, $\sim 3$min \\
\bottomrule
\end{tabular}
\end{table}

%----------------------------------------------------------------------
\subsection{Coordinate Descent for Group Multipliers}
\label{app:coord_descent}
%----------------------------------------------------------------------

The hyperparameter vector $(\lambda_{\mathrm{stat}}, \mu_1^{\mathrm{group}}, \ldots, \mu_G^{\mathrm{group}})$ is $G + 1 = 9$ dimensional.
Exhaustive grid search with $k$ values per dimension requires $k^9$ evaluations; even with $k = 5$, this exceeds $10^6$ configurations, each requiring one full Theory-KRR fit.

We adopt cyclic coordinate descent over the group multipliers, following the algorithmic template of \citet{friedman2010regularization}.
The statistical penalty $\lambda_{\mathrm{stat}}$ is fixed by standard KRR cross-validation (a one-dimensional problem).
The remaining 8 group multipliers are optimized sequentially: for each group $g$, we evaluate 16 candidate values
\begin{equation}
\label{eq:cd_candidates}
\mu_g^{\mathrm{group}} \in \{0\} \cup \{10^{u} : u \in \mathrm{linspace}(-4, 1, 15)\},
\end{equation}
holding all other multipliers fixed, and retain the value that minimizes validation-set MSE.
Including zero in the candidate set allows the algorithm to shut off entire theory families when they do not improve prediction.

Each evaluation requires one Nystr\"{o}m KRR solve ($m \times m$ Cholesky, $m = 500$) plus a single Newton-like correction that incorporates the structural penalty gradient (see Section~\ref{app:nystrom}).
The 16 candidates within each group are evaluated in parallel across CPU cores, so the wall-clock cost per group equals a single evaluation.
One full cycle through $G = 8$ groups therefore costs $8 \times 1$ evaluations in wall-clock time; we run up to 3 cycles and stop early if no multiplier changes.

The total budget is at most $3 \times 8 \times 16 = 384$ evaluations, completing in under 5 seconds per rolling window on our 40-core workstation with Nystr\"{o}m features.
Coordinate descent has three advantages over random or grid search: (i)~it exploits the separable structure of the grouped penalty, (ii)~it automatically resolves substitution effects between correlated theory families, and (iii)~it scales linearly in the number of groups rather than exponentially.

%----------------------------------------------------------------------
\subsection{Parallelization Strategy}
\label{app:parallel}
%----------------------------------------------------------------------

We exploit three levels of parallelism:

\begin{enumerate}
\item \textbf{GPU acceleration.}
The kernel matrix computation~\eqref{eq:nystrom_knm}--\eqref{eq:nystrom_kmm} involves pairwise distance calculations and element-wise exponentials, which are trivially parallelizable on the GPU.
The pairwise distance computation for $n = 90{,}000$ points against $m = 500$ landmarks completes in under 2 seconds on the RTX A6000, compared to approximately 30 seconds on a single CPU core.
The subsequent $m \times m$ Cholesky factorization is performed on the CPU, as the $m = 500$ system is small enough that GPU transfer overhead dominates.

\item \textbf{Multi-threaded coordinate descent.}
Within each coordinate descent cycle, the 16 candidate values for a given group multiplier $\mu_g^{\mathrm{group}}$ are independent and evaluated in parallel across CPU threads.
With 40 threads, all 16 candidates complete in the wall-clock time of a single Nystr\"{o}m solve.
The sequential component---cycling through the $G = 8$ groups---is inherently serial, but each group's line search is fully parallel, yielding an effective speedup of $\min(16, 40) = 16\times$ relative to sequential evaluation.

\item \textbf{Parallel ablation studies.}
For predefined $\mu$ configurations (e.g., single-group activations used in the grid search baseline), each configuration is independent.
We distribute all configurations across $\min(C, 40)$ CPU threads, achieving near-linear speedup since they share no mutable state.
\end{enumerate}

The memory footprint scales linearly in the number of parallel workers, since each worker maintains its own copy of the Nystr\"{o}m feature matrix $\mathbf{Z}$.
With $m = 500$ and $n = 90{,}000$, each copy occupies approximately 360~MB, so 40 parallel workers require 14~GB---well within the 274~GB available on our workstation.


\end{document}
